\chapter{Fungsi Kontinu} \label{lim:chapter}

%%%%%%%%%%%%%%%%%%%%%%%%%%%%%%%%%%%%%%%%%%%%%%%%%%%%%%%%%%%%%%%%%%%%%%%%%%%%%%

\section{Limit fungsi}
\label{sec:limoffunc}

%mbxINTROSUBSECTION

\sectionnotes{2--3 kuliah}

Sebelum mendefinisikan kekontinuan fungsi, kita akan membahas suatu gagasan
limit yang sedikit lebih umum daripada limit barisan.
Diberikan fungsi $f \colon S \to
\R$, kita ingin melihat perilaku $f(x)$ ketika $x$ mendekati suatu titik tertentu.

\subsection{Titik akumulasi}

Pertama,
kita kembali ke suatu konsep yang sebelumnya telah dijumpai dalam sebuah latihan.
Ketika bergerak di dalam himpunan $S$, kita hanya dapat mendekati titik-titik
yang di sekitarnya terdapat anggota $S$ sedekat apa pun.

\begin{defn}
Misalkan $S \subset \R$ suatu himpunan.  Bilangan $x \in \R$ disebut
\emph{\myindex{titik akumulasi}} dari $S$
jika untuk setiap $\epsilon > 0$, himpunan $(x-\epsilon,x+\epsilon) \cap S
\setminus \{ x \}$ tidak kosong.
\end{defn}

Artinya, $x$ merupakan titik akumulasi dari $S$ jika terdapat titik-titik dalam $S$
yang sedekat apa pun dengan $x$.  Definisi ini juga dapat dinyatakan dengan mengatakan
bahwa $x$ merupakan titik akumulasi dari $S$ jika untuk setiap $\epsilon > 0$, terdapat
$y \in S$ sedemikian sehingga $y \neq x$ dan $\sabs{x - y} < \epsilon$.
Perhatikan bahwa titik akumulasi dari $S$ tidak harus berada dalam $S$.

Mari kita lihat beberapa contoh.
\begin{enumerate}[(i)]
\item Himpunan
$\{ \nicefrac{1}{n} : n \in \N \}$ memiliki nol sebagai satu-satunya titik akumulasi.
\item Titik-titik akumulasi dari interval terbuka $(0,1)$ adalah
semua titik dalam interval tertutup $[0,1]$.
\item
Himpunan titik akumulasi dari $\Q$ adalah seluruh garis bilangan real $\R$.
\item 
Himpunan titik akumulasi dari $[0,1) \cup \{ 2 \}$ adalah interval $[0,1]$.
\item Himpunan $\N$ tidak memiliki titik akumulasi dalam $\R$.
\end{enumerate}

\begin{prop}
Misalkan $S \subset \R$.  Maka $x \in \R$ merupakan titik akumulasi dari $S$
jika dan hanya jika
terdapat barisan bilangan konvergen $\{ x_n \}_{n=1}^\infty$ sedemikian sehingga
$x_n \neq x$ dan $x_n \in S$ untuk semua $n$, serta $\lim\limits_{n\to\infty} x_n = x$.
\end{prop}

\begin{proof}
Pertama, andaikan $x$ merupakan titik akumulasi dari $S$.
Untuk setiap $n \in \N$, pilih $x_n$ sebagai sebarang titik dalam
$(x-\nicefrac{1}{n},x+\nicefrac{1}{n}) \cap S \setminus \{x\}$, yang
tidak kosong karena $x$ merupakan titik akumulasi dari $S$.
Maka
$x_n$ berada dalam jarak $\nicefrac{1}{n}$ dari $x$, yaitu,
\begin{equation*}
\sabs{x-x_n} < \nicefrac{1}{n} .
\avoidbreak
\end{equation*}
Karena $\{ \nicefrac{1}{n} \}_{n=1}^\infty$ konvergen ke nol,
$\{ x_n \}_{n=1}^\infty$ konvergen ke $x$.

Sebaliknya, jika kita mulai dengan suatu barisan bilangan
$\{ x_n \}_{n=1}^\infty$ dalam $S$
yang konvergen ke $x$ dan memenuhi $x_n \neq x$ untuk semua $n$, maka untuk setiap
$\epsilon > 0$ terdapat $M$ sedemikian sehingga, khususnya, $\sabs{x_M - x} <
\epsilon$.  Artinya, $x_M \in (x-\epsilon,x+\epsilon) \cap S \setminus \{x\}$.
\end{proof}

\subsection{Limit fungsi}

Jika fungsi $f$ didefinisikan pada suatu himpunan $S$ dan $c$ merupakan titik akumulasi dari $S$,
kita mendefinisikan limit $f(x)$ ketika $x$ mendekati $c$.
Apakah $f$ didefinisikan pada $c$ atau tidak, hal itu tidak berpengaruh pada definisi.
Bahkan jika fungsi tersebut didefinisikan pada $c$, limit
fungsi ketika $x$ menuju $c$ dapat saja berbeda
dari $f(c)$.

\begin{defn}
\index{limit!fungsi}%
Misalkan $f \colon S \to \R$ suatu fungsi dan $c$ suatu titik akumulasi dari
$S \subset \R$.
Andaikan terdapat $L \in \R$ dan untuk setiap $\epsilon > 0$,
terdapat $\delta > 0$ sedemikian sehingga setiap kali $x \in S \setminus \{ c \}$
dan $\sabs{x - c} < \delta$, berlaku
\begin{equation*}
\babs{f(x) - L} < \epsilon .
\end{equation*}
Kita kemudian mengatakan bahwa $f(x)$ \emph{konvergen}\index{konvergen!fungsi} ke $L$ ketika $x$
menuju $c$, dan menuliskan
\glsadd{not:limitasarrows}%
\begin{equation*}
f(x) \to L \quad\text{ketika}\quad x \to c .
\end{equation*}
Kita mengatakan bahwa $L$ merupakan \emph{limit} dari $f(x)$ ketika $x$
menuju $c$, dan jika $L$ tunggal (memang demikian), kita menuliskan
\glsadd{not:limfunc}%
\begin{equation*}
\lim_{x \to c} f(x) \coloneqq L .
\end{equation*}
Jika tidak terdapat $L$ semacam itu, kita mengatakan bahwa limitnya tidak ada atau
bahwa $f$ \emph{\myindex{divergen}} pada $c$.
\end{defn}

Sekali lagi, notasi dan bahasa yang kita gunakan di atas mengasumsikan bahwa limit $L$, jika
ada, bersifat tunggal; hal ini perlu dibuktikan.  Perhatikan bahwa fakta bahwa $c$
merupakan titik akumulasi penting untuk membuktikan ketunggalan.

\begin{prop}
Misalkan $c$ suatu titik akumulasi dari $S \subset \R$ dan misalkan $f \colon S \to \R$
suatu fungsi sedemikian sehingga $f(x)$ konvergen ketika $x$ menuju $c$.  Maka
limit $f(x)$ ketika $x$ menuju $c$ bersifat tunggal.
\end{prop}

\begin{proof}
Misalkan $L_1$ dan $L_2$ dua bilangan yang sama-sama memenuhi definisi.
Ambil sebarang $\epsilon > 0$ dan pilih $\delta_1 > 0$ sedemikian sehingga
$\babs{f(x)-L_1} < \nicefrac{\epsilon}{2}$ 
untuk semua $x \in S \setminus \{c\}$ dengan $\sabs{x-c} < \delta_1$.
Pilih pula $\delta_2 > 0$ sedemikian sehingga
$\babs{f(x)-L_2} < \nicefrac{\epsilon}{2}$
untuk semua $x \in S \setminus \{c\}$ dengan $\sabs{x-c} < \delta_2$.
Tetapkan $\delta \coloneqq \min \{ \delta_1, \delta_2 \}$.  Andaikan $x \in S$,
$\sabs{x-c} < \delta$, dan $x \neq c$.  Karena $\delta > 0$ dan $c$ merupakan titik
akumulasi, $x$ semacam itu ada.  Maka
\begin{equation*}
\sabs{L_1 - L_2} =
\babs{L_1 - f(x) + f(x) - L_2} \leq
\babs{L_1 - f(x)} + \babs{f(x) - L_2} < \frac{\epsilon}{2} + \frac{\epsilon}{2}
= \epsilon.
\end{equation*}
Karena $\sabs{L_1-L_2} < \epsilon$ untuk sebarang $\epsilon > 0$, kita peroleh
$L_1 = L_2$.
\end{proof}

\begin{example}
Perhatikan $f \colon \R \to \R$ yang didefinisikan oleh $f(x) \coloneqq x^2$.  Maka
untuk setiap $c \in \R$,
\begin{equation*}
\lim_{x\to c} f(x) = \lim_{x\to c} x^2 = c^2 .
\end{equation*}

Bukti: Tetapkan $c \in \R$, dan ambil sebarang $\epsilon > 0$.  Tuliskan
\begin{equation*}
\delta \coloneqq \min \left\{ 1 , \, \frac{\epsilon}{2\sabs{c}+1} \right\} .
\end{equation*}
Ambil $x \neq c$ sedemikian sehingga $\sabs{x-c} < \delta$.  Khususnya,
$\sabs{x-c} < 1$.  Berdasarkan ketaksamaan segitiga terbalik,
\begin{equation*}
\sabs{x}-\sabs{c} \leq \sabs{x-c} < 1 .
\end{equation*}
Dengan menambahkan $2\sabs{c}$ pada kedua ruas, kita peroleh
$\sabs{x} + \sabs{c} < 2\sabs{c} + 1$.  Taksir
\begin{equation*}
\begin{split}
\babs{f(x) - c^2} &= \sabs{x^2-c^2} \\
&= \babs{(x+c)(x-c)} \\
&= \sabs{x+c}\sabs{x-c} \\
&\leq \bigl(\sabs{x}+\sabs{c}\bigr)\sabs{x-c} \\
&< \bigl(2\sabs{c}+1\bigr)\sabs{x-c} \\
&< \bigl(2\sabs{c}+1\bigr)\frac{\epsilon}{2\sabs{c}+1} = \epsilon .
\end{split}
\end{equation*}
\end{example}

\begin{example}
Definisikan $f \colon [0,1) \to \R$ dengan
\begin{equation*}
f(x) \coloneqq 
\begin{cases}
x & \text{jika } x > 0 , \\
1 & \text{jika } x = 0 .
\end{cases}
\end{equation*}
Maka
$\lim\limits_{x\to 0} f(x) = 0$,
meskipun $f(0) = 1$.  Lihat \figureref{fig:limvaldiff}.
\begin{myfigureht}
\myincludegraphics{limvaldiff}{%
Grafik suatu fungsi yang bernilai 1 ketika variabel bebas
x bernilai 0; ketika variabel bebas
x lebih besar dari 0, fungsi tersebut sama dengan x.
Ditunjukkan bahwa fungsi tidak didefinisikan pada
x yang sama dengan 1 atau lebih besar dari 1, maupun pada x yang kurang dari 0.}
\caption{Fungsi dengan limit dan nilai yang berbeda pada $0$.\label{fig:limvaldiff}}
\end{myfigureht}

Bukti:  Ambil sebarang $\epsilon > 0$.  Misalkan $\delta \coloneqq \epsilon$.
Untuk $x \in [0,1)$, $x \neq 0$, dan $\sabs{x-0} < \delta$, kita peroleh
\begin{equation*}
\babs{f(x) - 0} = \sabs{x} < \delta = \epsilon .
\end{equation*}
\end{example}

\subsection{Limit sekuensial} \label{subseq:sequentiallimits}

Mari kita hubungkan limit sebagaimana didefinisikan di atas dengan limit barisan.

\begin{lemma}\label{seqflimit:lemma}
Misalkan $S \subset \R$, misalkan $c$ titik akumulasi dari $S$, misalkan $f \colon S \to
\R$ suatu fungsi, dan misalkan $L \in \R$.

Maka
$f(x) \to L$ ketika $x \to c$ jika dan hanya jika untuk setiap barisan
$\{ x_n \}_{n=1}^\infty$
sedemikian sehingga $x_n \in S \setminus \{c\}$ untuk semua $n$,
dan sedemikian sehingga $\lim_{n\to\infty} x_n = c$,
barisan $\bigl\{ f(x_n) \bigr\}_{n=1}^\infty$ konvergen ke $L$.
\end{lemma}

\begin{proof}
Misalkan
$f(x) \to L$ ketika $x \to c$, dan $\{ x_n \}_{n=1}^\infty$ suatu barisan
sedemikian sehingga
$x_n \in S \setminus \{c\}$ dan
$\lim_{n\to\infty} x_n = c$.
Kita hendak menunjukkan bahwa $\bigl\{ f(x_n) \bigr\}_{n=1}^\infty$ konvergen ke $L$.
Ambil $\epsilon > 0$.  Cari $\delta > 0$ sedemikian sehingga
jika $x \in S \setminus \{c\}$ dan $\sabs{x-c} < \delta$, maka
$\babs{f(x) - L} < \epsilon$.  Karena
$\{ x_n \}_{n=1}^\infty$ konvergen ke $c$, cari $M$ sedemikian sehingga untuk $n \geq M$,
berlaku $\sabs{x_n - c} < \delta$.  Oleh karena itu, untuk $n \geq M$,
\begin{equation*}
\babs{f(x_n) - L} < \epsilon .
\end{equation*}
Dengan demikian, $\bigl\{ f(x_n) \bigr\}_{n=1}^\infty$ konvergen ke $L$.

Untuk arah sebaliknya, kita menggunakan pembuktian dengan kontraposisi.  Misalkan
pernyataan bahwa $f(x) \to L$ ketika $x \to c$ tidak benar.  Negasi dari
definisi tersebut menyatakan bahwa terdapat $\epsilon > 0$ sedemikian sehingga untuk setiap
$\delta > 0$ terdapat $x \in S \setminus \{c\}$ dengan
$\sabs{x-c} < \delta$
dan $\babs{f(x)-L} \geq \epsilon$.

Gunakan $\nicefrac{1}{n}$ sebagai $\delta$ dalam pernyataan di atas untuk
membangun barisan $\{ x_n \}_{n=1}^\infty$.  Kita memperoleh
bahwa terdapat $\epsilon > 0$ sedemikian sehingga untuk setiap $n$,
terdapat titik $x_n \in S \setminus \{c\}$ dengan
$\sabs{x_n-c} < \nicefrac{1}{n}$
dan $\babs{f(x_n)-L} \geq \epsilon$.
Barisan $\{ x_n \}_{n=1}^\infty$ yang baru dibangun konvergen ke $c$, tetapi
barisan $\bigl\{ f(x_n) \bigr\}_{n=1}^\infty$ tidak konvergen ke $L$.
Dengan demikian, bukti selesai.
\end{proof}

Arah sebaliknya dalam lemma dapat diperkuat dengan cukup menyatakan bahwa
\myquote{$\bigl\{ f(x_n) \bigr\}_{n=1}^\infty$ konvergen,} tanpa mensyaratkan limit tertentu.
Lihat \exerciseref{exercise:seqflimitalt}.

\begin{example}
$\displaystyle \lim_{x \to 0} \sin( \nicefrac{1}{x} )$
tidak ada, tetapi
$\displaystyle \lim_{x \to 0} x\sin( \nicefrac{1}{x} ) = 0$.
Lihat \figureref{figsin1x}.

\begin{myfigureht}
%left guy also used in 10.4
\myincludepdft{sin1x_xsin1x}{%
Grafik fungsi-fungsi pada suatu interval di sekitar titik asal.
Di sebelah kiri, fungsi berosilasi
antara minus 1 dan 1
semakin cepat ketika x mendekati 0, sedemikian cepat sehingga di dekat titik asal
grafiknya tampak sebagai daerah hitam pekat.
Di sebelah kanan, fungsi kembali berosilasi semakin cepat
ketika x mendekati 0, tetapi kini amplitudonya juga menurun ke 0,
sehingga daerah yang sepintas tampak hitam pekat sebenarnya menyerupai
ujung segitiga.}
\caption{Grafik $\sin(\nicefrac{1}{x})$ dan $x \sin(\nicefrac{1}{x})$.
Perhatikan bahwa komputer tidak dapat menggambar $\sin(\nicefrac{1}{x})$ dengan tepat
di dekat nol karena fungsi tersebut berosilasi terlalu cepat.\label{figsin1x}}
\end{myfigureht}

Bukti:
Kita mulai dengan $\sin(\nicefrac{1}{x})$.  Definisikan suatu barisan
dengan
$x_n \coloneqq \frac{1}{\pi n + \nicefrac{\pi}{2}}$.  Mudah dilihat
bahwa $\lim_{n\to\infty} x_n = 0$.  Selanjutnya,
\begin{equation*}
\sin ( \nicefrac{1}{x_n} )
=
\sin (\pi n + \nicefrac{\pi}{2})
= {(-1)}^n .
\end{equation*}
Oleh karena itu, $\bigl\{ \sin ( \nicefrac{1}{x_n} ) \bigr\}_{n=1}^\infty$ tidak konvergen.
Berdasarkan
\lemmaref{seqflimit:lemma},
$\displaystyle \lim_{x \to 0} \sin( \nicefrac{1}{x} )$ tidak ada.

Sekarang tinjau $x\sin(\nicefrac{1}{x})$.  Misalkan $\{ x_n \}_{n=1}^\infty$ suatu barisan
sedemikian sehingga $x_n \neq 0$ untuk semua $n$, dan $\lim_{n\to\infty} x_n = 0$.
Perhatikan bahwa $\babs{\sin(t)} \leq 1$ untuk semua $t \in \R$.  Oleh karena itu,
\begin{equation*}
\babs{x_n\sin(\nicefrac{1}{x_n})-0}
=
\sabs{x_n}\babs{\sin(\nicefrac{1}{x_n})}
\leq
\sabs{x_n} .
\end{equation*}
Karena $x_n$ menuju nol, $\sabs{x_n}$ menuju nol, sehingga
$\bigl\{ x_n\sin(\nicefrac{1}{x_n}) \bigr\}_{n=1}^\infty$ konvergen ke nol.  Berdasarkan
\lemmaref{seqflimit:lemma},
$\displaystyle \lim_{x \to 0} x\sin( \nicefrac{1}{x} ) = 0$.
\end{example}

Ingatlah frasa \myquote{untuk setiap barisan} dalam lemma.
Sebagai contoh, ambil $\sin(\nicefrac{1}{x})$ dan barisan yang diberikan oleh
$x_n \coloneqq \nicefrac{1}{\pi n}$.
Maka $\bigl\{ \sin (\nicefrac{1}{x_n}) \bigr\}_{n=1}^\infty$
merupakan barisan nol konstan, sehingga
konvergen ke nol, tetapi limit
$\sin(\nicefrac{1}{x})$ ketika $x \to 0$ tidak ada.

Dengan menggunakan \lemmaref{seqflimit:lemma},
kita dapat mulai menerapkan semua yang kita ketahui tentang
limit barisan pada limit fungsi.  Mari kita berikan beberapa
contoh penting.

\begin{cor}
Misalkan $S \subset \R$ dan misalkan $c$ titik akumulasi dari $S$.
Misalkan $f \colon S \to
\R$ dan $g \colon S \to \R$ merupakan fungsi-fungsi
sedemikian sehingga
limit $f(x)$ dan $g(x)$ ketika $x$ menuju $c$ keduanya ada,
dan
\begin{equation*}
f(x) \leq g(x) \qquad \text{untuk semua } x \in S \setminus \{ c \}.
\end{equation*}
Maka
\begin{equation*}
\lim_{x\to c} f(x) \leq \lim_{x\to c} g(x) .
\end{equation*}
\end{cor}

\begin{proof}
Ambil $\{ x_n \}_{n=1}^\infty$ sebagai barisan bilangan dalam $S \setminus \{ c \}$
yang konvergen ke $c$.  Misalkan
\begin{equation*}
L_1 \coloneqq \lim_{x\to c} f(x), \qquad \text{dan} \qquad L_2 \coloneqq \lim_{x\to c} g(x) .
\end{equation*}
\lemmaref{seqflimit:lemma} menyatakan bahwa $\bigl\{ f(x_n) \bigr\}_{n=1}^\infty$ konvergen ke
$L_1$ dan $\bigl\{ g(x_n) \bigr\}_{n=1}^\infty$ konvergen ke $L_2$.  Kita juga
mempunyai $f(x_n) \leq g(x_n)$ untuk semua $n$.
Kita memperoleh $L_1 \leq L_2$ dengan menggunakan
\lemmaref{limandineq:lemma}.
\end{proof}

Dengan menerapkan fungsi konstan, kita memperoleh korolari berikut.  Buktinya
diserahkan sebagai latihan.

\begin{cor} \label{fconstineq:cor}
Misalkan $S \subset \R$ dan misalkan $c$ titik akumulasi dari $S$.  Misalkan $f \colon S \to
\R$ suatu fungsi sedemikian sehingga limit $f(x)$ ketika $x$ menuju $c$
ada.
Misalkan terdapat dua bilangan real $a$ dan $b$ sedemikian sehingga
\begin{equation*}
a \leq f(x) \leq b \qquad \text{untuk semua } x \in S \setminus \{ c \}.
\end{equation*}
Maka
\begin{equation*}
a \leq \lim_{x\to c} f(x) \leq b .
\end{equation*}
\end{cor}

Dengan menggunakan \lemmaref{seqflimit:lemma} dengan cara yang sama seperti di atas, kita juga memperoleh
korolari-korolari berikut, yang buktinya kembali diserahkan sebagai latihan.

\begin{cor} \label{fsqueeze:cor}
Misalkan $S \subset \R$ dan misalkan $c$ titik akumulasi dari $S$.
Misalkan $f \colon S \to \R$,
$g \colon S \to \R$, dan $h \colon S \to \R$ merupakan fungsi-fungsi sedemikian sehingga
\begin{equation*}
f(x) \leq g(x) \leq h(x) \qquad \text{untuk semua } x \in S \setminus \{ c \}.
\end{equation*}
Misalkan limit $f(x)$ dan $h(x)$ ketika $x$ menuju $c$ keduanya ada, dan
\begin{equation*}
\lim_{x\to c} f(x) = \lim_{x\to c} h(x) .
\end{equation*}
Maka limit $g(x)$ ketika $x$ menuju $c$ ada dan
\begin{equation*}
\lim_{x\to c} g(x) =
\lim_{x\to c} f(x) = \lim_{x\to c} h(x) .
\end{equation*}
\end{cor}

\begin{cor} \label{falg:cor}
Misalkan $S \subset \R$ dan misalkan $c$ titik akumulasi dari $S$.
Misalkan $f \colon S \to \R$ dan
$g \colon S \to \R$ merupakan fungsi-fungsi
sedemikian sehingga
limit $f(x)$ dan $g(x)$ ketika $x$ menuju $c$ keduanya ada.
Maka
\begin{enumerate}[(i)]
\item
$\displaystyle
\lim_{x\to c} \bigl(f(x)+g(x)\bigr) = \left(\lim_{x\to c} f(x)\right) + 
\left(\lim_{x\to c} g(x)\right) .
$
\item
$\displaystyle
\lim_{x\to c} \bigl(f(x)-g(x)\bigr) = \left(\lim_{x\to c} f(x)\right) -
\left(\lim_{x\to c} g(x)\right) .
$
\item
$\displaystyle
\lim_{x\to c} \bigl(f(x)g(x)\bigr) = \left(\lim_{x\to c} f(x)\right)
\left(\lim_{x\to c} g(x)\right) .
$
\item \label{falg:cor:iv} Jika
$\displaystyle \lim_{x\to c} g(x) \neq 0$
dan $g(x) \neq 0$ untuk semua $x \in S \setminus \{ c \}$, maka
\begin{equation*}
\lim_{x\to c} \frac{f(x)}{g(x)} =
\frac{\lim_{x\to c} f(x)}{\lim_{x\to c} g(x)} .
\end{equation*}
\end{enumerate}
\end{cor}

\begin{cor} \label{fabs:cor}
Misalkan $S \subset \R$ dan misalkan $c$ titik akumulasi dari $S$.
Misalkan $f \colon S \to \R$ suatu fungsi sedemikian sehingga limit $f(x)$ ketika $x$ menuju $c$
ada.
Maka
\begin{equation*}
\lim_{x\to c} \babs{f(x)} =
\abs{\lim_{x\to c} f(x)}.
\end{equation*}
\end{cor}

\subsection{Limit pembatasan dan limit sepihak}

Kadang-kadang kita bekerja dengan fungsi yang didefinisikan pada suatu subhimpunan.

\begin{defn}
Misalkan $f \colon S \to \R$ suatu fungsi dan $A \subset S$.  Definisikan
fungsi $f|_A \colon A \to \R$ dengan
\begin{equation*}
f|_A (x) \coloneqq f(x)  \qquad \text{untuk } x \in A.
\end{equation*}
Kita menyebut $f|_A$ sebagai \emph{\myindex{pembatasan}} $f$ pada $A$.
\end{defn}

Fungsi $f|_A$ hanyalah fungsi $f$ yang diambil pada domain yang lebih kecil.
Proposisi berikut merupakan analog dari pengambilan ekor suatu barisan.
Proposisi tersebut mengatakan bahwa limit bersifat \myquote{lokal}: limit hanya bergantung pada titik-titik
yang sedekat apa pun dengan $c$.

\begin{prop} \label{prop:limrest}
Misalkan $S \subset \R$, $c \in \R$, dan
misalkan $f \colon S
\to \R$ suatu fungsi.
Andaikan
$A \subset S$ sedemikian sehingga untuk suatu $\alpha > 0$, berlaku
$(A \setminus \{ c \}) \cap (c-\alpha,c+\alpha) = (S \setminus \{ c \}) \cap (c-\alpha,c+\alpha)$.
\begin{enumerate}[(i)]
\item Titik $c$ merupakan titik akumulasi dari $A$ jika dan hanya jika $c$ merupakan titik akumulasi
dari $S$.
\item Jika $c$ merupakan titik akumulasi dari $S$, maka $f(x) \to L$ ketika $x \to c$ jika dan hanya jika
$f|_A(x) \to L$ ketika $x \to c$.
\end{enumerate}
\end{prop}

\begin{proof}
Pertama, misalkan $c$ suatu titik akumulasi dari $A$.
Karena $A \subset S$, jika $( A \setminus \{ c\} ) \cap
(c-\epsilon,c+\epsilon)$ tidak kosong untuk setiap $\epsilon > 0$,
maka $( S \setminus \{ c\} ) \cap
(c-\epsilon,c+\epsilon)$ tidak kosong untuk setiap $\epsilon > 0$.
Jadi, $c$ merupakan titik akumulasi dari $S$.
Kedua, andaikan $c$ suatu titik akumulasi
dari $S$.  Maka untuk $\epsilon > 0$ sedemikian sehingga $\epsilon < \alpha$,
kita peroleh $( A \setminus \{ c\} ) \cap (c-\epsilon,c+\epsilon) =
( S \setminus \{ c\} ) \cap (c-\epsilon,c+\epsilon)$, yang tidak kosong.  Hal ini berlaku untuk semua
$\epsilon < \alpha$ sehingga
$( A \setminus \{ c\} ) \cap (c-\epsilon,c+\epsilon)$ pasti tidak kosong untuk semua
$\epsilon > 0$.  Jadi, $c$ merupakan titik akumulasi dari $A$.

Sekarang andaikan $c$ suatu titik akumulasi dari $S$ dan $f(x) \to L$ ketika $x \to c$.  Artinya, untuk setiap $\epsilon > 0$
terdapat $\delta > 0$ sedemikian sehingga jika $x \in S \setminus \{ c \}$
dan $\sabs{x-c} < \delta$, maka $\babs{f(x)-L} < \epsilon$.  Karena $A \subset S$,
jika $x \in A \setminus \{ c \}$, maka $x \in S \setminus \{ c \}$,
sehingga $f|_A(x) \to L$ ketika $x \to c$.

Terakhir, andaikan $f|_A(x) \to L$ ketika $x \to c$ dan ambil sebarang
$\epsilon > 0$.
Terdapat $\delta' > 0$ sedemikian sehingga jika $x \in A \setminus \{ c \}$
dan $\sabs{x-c} < \delta'$, maka $\babs{ f|_A(x)-L } < \epsilon$.
Ambil $\delta \coloneqq \min \{ \delta', \alpha \}$.
Sekarang andaikan $x \in S \setminus \{ c \}$ dan
$\sabs{x-c} < \delta$.  Karena $\sabs{x-c} < \alpha$, kita peroleh $x \in A \setminus \{ c \}$,
dan karena $\sabs{x-c} < \delta'$,
kita peroleh $\babs{f(x)-L} = \babs{ f|_A(x)-L } < \epsilon$.
\end{proof}

Hipotesis mengenai $A$ dalam proposisi tersebut diperlukan.  Untuk pembatasan
sebarang, umumnya kita hanya memperoleh implikasi dalam satu arah
(lihat \exerciseref{exercise:restrictionlimitexercise}).
Notasi yang lazim untuk limit tersebut adalah
\begin{equation*}
\lim_{\substack{x \to c\\x \in A}} f(x) \coloneqq \lim_{x \to c} f|_A(x) .
\end{equation*}
Penggunaan pembatasan yang umum dalam konteks limit, yang tidak memenuhi
hipotesis dalam proposisi tersebut, adalah
\emph{\myindex{limit sepihak}}%
\footnote{%
Terdapat beraneka ragam notasi limit sepihak.
Sebagai contoh,
$\lim\limits_{\substack{x \to c\\x < c}} f(x)$,
$\lim\limits_{x \uparrow c} f(x)$, atau
$\lim\limits_{x \nearrow c} f(x)$ untuk
$\lim\limits_{x \to c^-} f(x)$.}.

\begin{defn} \label{defn:onesidedlimits}
Misalkan $f \colon S \to \R$ suatu fungsi dan misalkan $c \in \R$.  Jika
$c$ merupakan titik akumulasi dari
$S \cap (c,\infty)$
dan limit
pembatasan $f$ pada $S \cap (c,\infty)$
ketika $x \to c$ ada, definisikan
\glsadd{not:onesidedlim}
\begin{equation*}
\lim_{x \to c^+} f(x) \coloneqq \lim_{x\to c} f|_{S \cap (c,\infty)}(x) .
\end{equation*}
Demikian pula, jika $c$ merupakan titik akumulasi dari
$S \cap (-\infty,c)$ dan limit pembatasannya ketika $x \to c$
ada, definisikan
\begin{equation*}
\lim_{x \to c^-} f(x) \coloneqq \lim_{x\to c} f|_{S \cap (-\infty,c)}(x) .
\end{equation*}
\end{defn}

\propref{prop:limrest} tidak berlaku untuk limit sepihak.
Limit-limit sepihak dapat saja ada, tetapi limit pada suatu titik tidak ada.  Sebagai
contoh, definisikan $f \colon \R \to \R$ dengan $f(x) \coloneqq 1$ untuk $x < 0$ dan
$f(x) \coloneqq 0$ untuk $x \geq 0$.
Pembaca dipersilakan memverifikasi bahwa
\begin{equation*}
\lim_{x \to 0^-} f(x) = 1, \qquad
\lim_{x \to 0^+} f(x) = 0, \qquad
\lim_{x \to 0} f(x) \quad \text{tidak ada.}
\end{equation*}
Meskipun demikian, kita memiliki pengganti berikut.

\begin{prop} \label{prop:onesidedlimits}
Misalkan $S \subset \R$ sedemikian sehingga $c$ merupakan titik akumulasi
dari $S \cap (-\infty,c)$ dan $S \cap (c,\infty)$, misalkan
$f \colon S \to \R$ suatu fungsi, dan misalkan $L \in \R$.  Maka $c$ merupakan titik akumulasi dari $S$ dan
\begin{equation*}
\lim_{x \to c} f(x) = L
\qquad \text{jika dan hanya jika} \qquad
\lim_{x \to c^-} f(x) =
\lim_{x \to c^+} f(x) =
L .
\end{equation*}
\end{prop}

Artinya, limit pada $c$ ada jika dan hanya jika kedua limit sepihak ada dan bernilai sama.  Pembuktiannya
merupakan penerapan langsung dari definisi limit
dan diserahkan sebagai latihan.  Pokok utamanya adalah bahwa
$\bigl( S \cap (-\infty,c) \bigr) \cup \bigl( S \cap (c,\infty) \bigr)
= S \setminus \{ c \}$.

\subsection{Latihan}

\begin{exercise}
Tentukan limitnya (dan tentu saja buktikan) atau buktikan bahwa limit tersebut
tidak ada:

\medskip

\noindent
\begin{tabular}{lllll}
a)
$\displaystyle
\lim_{x\to c} \sqrt{x}
$, untuk $c \geq 0$
& &
b)
$\displaystyle
\lim_{x\to c} x^2+x+1
$, untuk $c \in \R$
& &
c)
$\displaystyle
\lim_{x\to 0} x^2 \cos (\nicefrac{1}{x})
$
\\
d)
$\displaystyle
\lim_{x\to 0} \sin(\nicefrac{1}{x}) \cos (\nicefrac{1}{x})
$
& &
e)
$\displaystyle
\lim_{x\to 0} \sin(x) \cos (\nicefrac{1}{x})
$ & 
\end{tabular}
\end{exercise}

\begin{exercise}
Buktikan \corref{fconstineq:cor}.
\end{exercise}

\begin{exercise}
Buktikan \corref{fsqueeze:cor}.
\end{exercise}

\begin{exercise}
Buktikan \corref{falg:cor}.
\end{exercise}

\begin{exercise}
Misalkan $A \subset S$.  Tunjukkan bahwa jika $c$ merupakan titik akumulasi dari $A$, maka $c$
merupakan titik akumulasi dari $S$.  Perhatikan perbedaannya dengan
\propref{prop:limrest}.
\end{exercise}

\begin{exercise} \label{exercise:restrictionlimitexercise}
Misalkan $A \subset S$.  Andaikan $c$ merupakan titik akumulasi dari $A$ dan
juga merupakan titik akumulasi dari $S$.
Misalkan $f \colon S \to \R$ suatu fungsi.  Tunjukkan bahwa jika
$f(x) \to L$ ketika $x \to c$, maka
$f|_A(x) \to L$ ketika $x \to c$.
Perhatikan perbedaannya dengan
\propref{prop:limrest}.
\end{exercise}

\begin{exercise}
Berikan contoh fungsi $f \colon [-1,1] \to \R$ sedemikian sehingga untuk
$A\coloneqq [0,1]$, berlaku
$f|_A(x) \to 0$ ketika $x \to 0$, tetapi limit $f(x)$ ketika $x \to 0$
tidak ada.  Jelaskan mengapa Anda tidak dapat menerapkan
\propref{prop:limrest}.
\end{exercise}

\begin{exercise}
Berikan contoh fungsi $f$ dan $g$ sedemikian sehingga limit $f(x)$
maupun $g(x)$ tidak ada ketika $x \to 0$, tetapi limit $f(x)+g(x)$ ada
ketika $x \to 0$.
\end{exercise}

\begin{exercise} \label{exercise:contlimitcomposition}
Misalkan $c_1$ merupakan titik akumulasi dari $A \subset \R$ dan $c_2$
merupakan titik akumulasi dari $B \subset \R$.  Andaikan
$f \colon A \to B$ dan $g \colon B \to \R$ merupakan fungsi
sedemikian sehingga
$f(x) \to c_2$ ketika $x \to c_1$ dan
$g(y) \to L$ ketika $y \to c_2$.  Jika $c_2 \in B$, andaikan juga bahwa $g(c_2) = L$.
Misalkan $h(x) \coloneqq g\bigl(f(x)\bigr)$ dan tunjukkan
$h(x) \to L$ ketika $x \to c_1$.
Petunjuk: Perhatikan bahwa $f(x)$ dapat sama dengan $c_2$ untuk banyak $x \in A$
(lihat juga
\exerciseref{exercise:contlimitbadcomposition}).
\end{exercise}

\begin{exercise}
Andaikan $f \colon \R \to \R$ merupakan fungsi sedemikian sehingga untuk setiap
barisan $\{x_n\}_{n=1}^\infty$ dalam $\R$, barisan
$\bigl\{ f(x_n) \bigr\}_{n=1}^\infty$ konvergen.  Buktikan bahwa
$f$ konstan; artinya,
$f(x) = f(y)$ untuk semua $x,y \in \R$.
\end{exercise}

\begin{exercise} \label{exercise:seqflimitalt}
Buktikan versi yang lebih kuat berikut dari salah satu arah dalam
\lemmaref{seqflimit:lemma}:
Misalkan $S \subset \R$, $c$ merupakan titik akumulasi dari $S$, dan $f \colon S \to
\R$ suatu fungsi.
Andaikan bahwa untuk setiap barisan $\{x_n\}_{n=1}^\infty$ dalam $S \setminus \{c\}$ sedemikian sehingga
$\lim_{n\to\infty} x_n = c$, barisan $\bigl\{ f(x_n) \bigr\}_{n=1}^\infty$ konvergen.
Tunjukkan bahwa limit $f(x)$ ketika $x \to c$ ada.
\end{exercise}

\begin{exercise}
Buktikan \propref{prop:onesidedlimits}.
\end{exercise}

\begin{exercise}
Andaikan $S \subset \R$ dan $c$ merupakan titik akumulasi dari $S$.  Andaikan $f \colon
S \to \R$ terbatas.  Tunjukkan bahwa terdapat barisan
$\{ x_n \}_{n=1}^\infty$
dengan $x_n \in S \setminus \{ c \}$ dan $\lim_{n\to\infty} x_n = c$ sedemikian sehingga
$\bigl\{ f(x_n) \bigr\}_{n=1}^\infty$ konvergen.
\end{exercise}

\begin{exercise}[Menantang] \label{exercise:contlimitbadcomposition}
Tunjukkan bahwa asumsi $g(c_2) = L$ dalam
\exerciseref{exercise:contlimitcomposition} diperlukan.  Artinya, berikan $f$
dan $g$ sedemikian sehingga $f(x) \to c_2$ ketika $x \to c_1$ dan
$g(y) \to L$ ketika $y \to c_2$, tetapi $g\bigl(f(x)\bigr)$ tidak menuju $L$
ketika $x \to c_1$.
\end{exercise}

\begin{exercise}
Tunjukkan bahwa syarat titik akumulasi diperlukan agar definisi limit
masuk akal.  Artinya, andaikan $c$ bukan titik akumulasi
dari $S \subset \R$, dan $f \colon S \to \R$ suatu fungsi.  Tunjukkan bahwa
setiap $L$ akan memenuhi definisi limit di $c$ jika syarat
bahwa $c$ merupakan titik akumulasi ditiadakan.
\end{exercise}

\begin{exercise}
\leavevmode
\begin{enumerate}[a)]
\item
Buktikan \corref{fabs:cor}.
\item
Berikan contoh yang menunjukkan bahwa kebalikan
korolari tersebut tidak berlaku.
\end{enumerate}
\end{exercise}

%%%%%%%%%%%%%%%%%%%%%%%%%%%%%%%%%%%%%%%%%%%%%%%%%%%%%%%%%%%%%%%%%%%%%%%%%%%%%%

\sectionnewpage
\section{Fungsi kontinu}
\label{sec:cont}

%mbxINTROSUBSECTION

\sectionnotes{2--2,5 kuliah}

Kriteria tingkat sekolah menengah untuk konsep kekontinuan
adalah bahwa suatu fungsi kontinu jika
kita dapat menggambar grafiknya tanpa mengangkat pena dari kertas.  Meskipun
konsep intuitif tersebut dapat berguna dalam situasi sederhana, kita memerlukan
ketelitian.  Perumusan definisi berikut secara tepat memerlukan kontribusi tiga matematikawan besar
(Bolzano, Cauchy, dan akhirnya Weierstrass),
dan bentuk akhirnya baru muncul pada akhir abad ke-19.

\subsection{Definisi dan sifat-sifat dasar}

\begin{defn}
Misalkan $S \subset \R$ dan $c \in S$.
Kita mengatakan $f \colon S \to \R$
\emph{kontinu di $c$}\index{kontinu di $c$}
jika untuk setiap $\epsilon > 0$
terdapat $\delta > 0$ sedemikian sehingga setiap kali $x \in S$ dan $\sabs{x-c} <
\delta$, berlaku
$\babs{f(x)-f(c)} < \epsilon$.

%\medskip

Jika $f \colon S \to \R$ kontinu di setiap $c \in S$, kita cukup mengatakan
$f$ merupakan \emph{\myindex{fungsi kontinu}}\index{fungsi!kontinu}.
\end{defn}
\begin{myfigureht}
\myincludepdft{contigr}{%
Diagram grafik suatu fungsi y sama dengan f dari x,
di dekat titik x sama dengan c.  Grafiknya ditampilkan, beserta
pita dalam jarak delta dari c dan pita dalam jarak
epsilon dari f dari c.  Kotak yang berada di dalam
kedua pita diarsir, dan dalam rentang ini grafik
berada di dalam rentang vertikal.}
\caption{Untuk $\sabs{x-c} < \delta$, grafik $f(x)$ harus berada di dalam daerah abu-abu.\label{fig:contigr}}
\end{myfigureht}

Jika $f$ kontinu untuk semua $c \in A$, kita mengatakan
$f$ \emph{kontinu pada $A \subset S$}.  Sebuah
latihan langsung (\exerciseref{exercise:restrictioncontinuous})
menunjukkan bahwa hal ini mengakibatkan $f|_A$ kontinu, meskipun
pernyataan sebaliknya tidak berlaku (seperti akan kita lihat dalam \exampleref{example:removablediscont}).

Kekontinuan mungkin merupakan definisi terpenting yang perlu dipahami dalam analisis,
dan definisi ini tidak mudah.  Lihat \figureref{fig:contigr}.
Perhatikan bahwa $\delta$ tidak hanya bergantung
pada $\epsilon$, tetapi juga pada $c$; kita tidak perlu memilih
satu $\delta$ untuk semua $c \in S$.
Bukan suatu kebetulan
bahwa definisi kekontinuan mirip dengan definisi
limit fungsi.  Ciri utama kekontinuan
adalah bahwa fungsi kontinu tepat merupakan fungsi yang berperilaku baik terhadap limit.

\begin{prop} \label{contbasic:prop}
Tinjau fungsi $f \colon S \to \R$ yang didefinisikan pada himpunan $S \subset \R$
dan misalkan $c \in S$.
Maka:
%\begin{enumerate}[(i),itemsep=0.5\itemsep,parsep=0.5\parsep,topsep=0.5\topsep,partopsep=0.5\partopsep]
\begin{enumerate}[(i)]
\item \label{contbasic:prop:i}
Jika $c$ bukan titik akumulasi dari $S$, maka $f$ kontinu di $c$.
\item \label{contbasic:prop:ii}
Jika $c$ titik akumulasi dari $S$, maka $f$ kontinu di $c$
jika dan hanya jika limit $f(x)$ ketika $x \to c$ ada dan
\begin{equation*}
\lim_{x\to c} f(x) = f(c) .
\end{equation*}
\item \label{contbasic:prop:iii}
Fungsi $f$ kontinu di $c$ jika dan hanya jika untuk setiap barisan
$\{ x_n \}_{n=1}^\infty$
dengan $x_n \in S$ dan $\lim\limits_{n\to\infty} x_n = c$, barisan
$\bigl\{ f(x_n) \bigr\}_{n=1}^\infty$ konvergen
ke $f(c)$.
\end{enumerate}
\end{prop}

\begin{proof}
\pagebreak[2]
Kita mulai dengan \ref{contbasic:prop:i}.  Misalkan $c$ bukan titik akumulasi dari
$S$.  Maka terdapat $\delta > 0$
sedemikian sehingga $S \cap (c-\delta,c+\delta) = \{ c \}$.
Untuk sebarang $\epsilon > 0$, cukup pilih $\delta$ tersebut.
Satu-satunya $x \in S$ sedemikian sehingga $\sabs{x-c} < \delta$ adalah $x=c$.  Maka
$\babs{f(x)-f(c)} = \babs{f(c)-f(c)} = 0 < \epsilon$.

Berikutnya, tinjau \ref{contbasic:prop:ii}.
Misalkan $c$ titik akumulasi dari $S$.  Mula-mula, misalkan
$\lim_{x\to c} f(x) = f(c)$.  Maka untuk setiap $\epsilon > 0$,
terdapat $\delta > 0$ sedemikian sehingga jika $x \in S \setminus \{ c \}$
dan $\sabs{x-c} < \delta$, maka $\babs{f(x)-f(c)} < \epsilon$.
Selain itu, $\babs{f(c)-f(c)} = 0 < \epsilon$, sehingga definisi kekontinuan di
$c$ terpenuhi.  Sebaliknya, misalkan $f$ kontinu
di $c$.  Untuk setiap $\epsilon > 0$, terdapat $\delta > 0$
sedemikian sehingga untuk $x \in S$ dengan $\sabs{x-c} < \delta$, berlaku
$\babs{f(x)-f(c)} < \epsilon$.  Tentu saja, pernyataan tersebut tetap benar jika
$x \in S \setminus \{ c \} \subset S$.  Oleh karena itu, $\lim_{x\to c} f(x) =
f(c)$.

Untuk \ref{contbasic:prop:iii}, mula-mula misalkan $f$ kontinu di $c$.
Misalkan $\{ x_n \}_{n=1}^\infty$
suatu barisan sedemikian sehingga $x_n \in S$ dan $\lim_{n\to\infty} x_n = c$.  Ambil $\epsilon > 0$.
Cari $\delta > 0$ sedemikian sehingga $\babs{f(x)-f(c)} < \epsilon$
untuk semua $x \in S$ dengan $\sabs{x-c} < \delta$.  Cari $M \in \N$
sedemikian sehingga untuk $n \geq M$, berlaku $\sabs{x_n-c} < \delta$.  Maka untuk
$n \geq M$, berlaku $\babs{f(x_n)-f(c)} < \epsilon$,
sehingga $\bigl\{ f(x_n) \bigr\}_{n=1}^\infty$
konvergen ke $f(c)$.

Arah sebaliknya dari \ref{contbasic:prop:iii} kita buktikan dengan kontraposisi.
Misalkan $f$ tidak
kontinu di $c$.  Maka terdapat $\epsilon > 0$
sedemikian sehingga untuk setiap $\delta > 0$, terdapat $x \in S$
sedemikian sehingga $\sabs{x-c} < \delta$ dan $\babs{f(x)-f(c)} \geq \epsilon$.
Definisikan barisan $\{ x_n \}_{n=1}^\infty$ sebagai berikut.
Ambil $x_n \in S$ sedemikian sehingga $\sabs{x_n-c} < \nicefrac{1}{n}$
dan $\babs{f(x_n)-f(c)} \geq \epsilon$.
Sekarang $\{ x_n \}_{n=1}^\infty$ merupakan
barisan dalam $S$ sedemikian sehingga
$\lim_{n\to\infty} x_n = c$ dan
$\babs{f(x_n)-f(c)} \geq \epsilon$ untuk semua $n \in \N$.
Dengan demikian, $\bigl\{ f(x_n) \bigr\}_{n=1}^\infty$
tidak konvergen ke $f(c)$.  Barisan tersebut mungkin konvergen atau mungkin tidak, tetapi pasti
tidak konvergen ke $f(c)$.
\end{proof}

Butir terakhir dalam proposisi tersebut sangat kuat.  Butir itu memungkinkan kita
menerapkan dengan cepat apa yang kita ketahui tentang limit barisan pada fungsi kontinu,
bahkan untuk membuktikan bahwa fungsi tertentu kontinu.
Butir tersebut juga dapat diperkuat (lihat \exerciseref{exercise:contseqalt}).

\begin{example}
Fungsi $f \colon (0,\infty) \to \R$
yang didefinisikan oleh $f(x) \coloneqq \nicefrac{1}{x}$ bersifat kontinu.

Bukti: Tetapkan $c \in (0,\infty)$.
Misalkan $\{ x_n \}_{n=1}^\infty$ suatu barisan dalam $(0,\infty)$ sedemikian sehingga
$\lim_{n\to\infty} x_n = c$.  Maka
\begin{equation*}
f(c) = \frac{1}{c}
=
\frac{1}{\lim_{n\to\infty} x_n}
=
\lim_{n \to \infty} \frac{1}{x_n}
=
\lim_{n \to \infty} f(x_n) .
\end{equation*}
Dengan demikian, $f$ kontinu di $c$.  Karena $f$ kontinu di semua $c \in
(0,\infty)$, $f$ merupakan fungsi kontinu.
\end{example}

Sebelumnya, kita telah menunjukkan $\lim_{x \to c} x^2 = c^2$ secara langsung.  Oleh karena itu,
fungsi $x^2$ kontinu.  Butir terakhir dari \propref{contbasic:prop} dan kekontinuan
operasi aljabar terhadap limit barisan,
\propref{prop:contalg}, memberikan pembuktian singkat bagi hasil yang jauh lebih umum.

\begin{prop}
Misalkan $f \colon \R \to \R$ suatu \emph{\myindex{polinom}}.  Artinya,
\begin{equation*}
f(x) = a_d x^d + a_{d-1} x^{d-1} + \cdots + a_1 x + a_0 ,
\end{equation*}
untuk beberapa konstanta $a_0, a_1, \ldots, a_d$.
Maka $f$ kontinu.
\end{prop}

\begin{proof}
Tetapkan $c \in \R$.
Misalkan $\{ x_n \}_{n=1}^\infty$ suatu barisan sedemikian sehingga
$\lim_{n\to\infty} x_n = c$.  Maka
\begin{equation*}
\begin{split}
f(c) &=
a_d c^d + a_{d-1} c^{d-1} + \cdots + a_1 c + a_0 
\\
&= 
a_d {\left(\lim_{n\to\infty} x_n\right)}^d
+ a_{d-1} {\left(\lim_{n\to\infty} x_n\right)}^{d-1}
+ \cdots
+ a_1 \left(\lim_{n\to\infty} x_n\right) + a_0 
\\
& =
\lim_{n \to \infty}
\left(
a_d x_n^d + a_{d-1} x_n^{d-1} + \cdots + a_1 x_n + a_0 
\right)
=
\lim_{n \to \infty}
f(x_n) .
\avoidbreak
\end{split}
\end{equation*}
Dengan demikian, $f$ kontinu di $c$.  Karena $f$ kontinu di semua $c \in \R$,
$f$ merupakan fungsi kontinu.
\end{proof}

Dengan penalaran serupa, atau dengan menggunakan \corref{falg:cor},
kita dapat membuktikan proposisi berikut.  Buktinya diserahkan sebagai
latihan.

\begin{prop} \label{contalg:prop}
Misalkan $f \colon S \to \R$ dan $g \colon S \to \R$ fungsi-fungsi yang
kontinu di $c \in S$.
\begin{enumerate}[(i)]
\item Fungsi $h \colon S \to \R$ yang didefinisikan oleh
$h(x) \coloneqq f(x)+g(x)$ kontinu di $c$.
\item Fungsi $h \colon S \to \R$ yang didefinisikan oleh
$h(x) \coloneqq f(x)-g(x)$ kontinu di $c$.
\item Fungsi $h \colon S \to \R$ yang didefinisikan oleh
$h(x) \coloneqq f(x)g(x)$ kontinu di $c$.
\item Jika $g(x) \neq 0$ untuk semua $x \in S$, fungsi $h \colon S \to \R$
yang diberikan oleh $h(x) \coloneqq \frac{f(x)}{g(x)}$ kontinu di $c$.
\end{enumerate}
\end{prop}

\begin{example} \label{sincos:example}
Fungsi $\sin(x)$ dan $\cos(x)$ kontinu.
Dalam perhitungan berikut, kita menggunakan identitas trigonometri
yang mengubah jumlah menjadi hasil kali.  Kita juga menggunakan fakta sederhana bahwa
$\babs{\sin(x)} \leq \sabs{x}$, $\babs{\cos(x)} \leq 1$,
dan $\babs{\sin(x)} \leq 1$.
\begin{equation*}
\begin{split}
\babs{\sin(x)-\sin(c)} & =
\abs{
2 \sin \left( \frac{x-c}{2} \right) \cos \left( \frac{x+c}{2} \right)
}
\\
& =
2
\abs{ \sin \left( \frac{x-c}{2} \right) }
\abs{ \cos \left( \frac{x+c}{2} \right) }
\\
& \leq
2
\abs{ \sin \left( \frac{x-c}{2} \right) }
\\
& \leq
2
\abs{ \frac{x-c}{2} }
= \sabs{x-c}
\end{split}
\end{equation*}
\begin{equation*}
\begin{split}
\babs{\cos(x)-\cos(c)} & =
\abs{
-2 \sin \left( \frac{x-c}{2} \right) \sin \left( \frac{x+c}{2} \right)
}
\\
& =
2
\abs{ \sin \left( \frac{x-c}{2} \right) }
\abs{ \sin \left( \frac{x+c}{2} \right) }
\\
& \leq
2
\abs{ \sin \left( \frac{x-c}{2} \right) }
\\
& \leq
2
\abs{ \frac{x-c}{2} }
= \sabs{x-c}
\end{split}
\end{equation*}

Pernyataan bahwa $\sin$ dan $\cos$ kontinu diperoleh dengan mengambil
sebarang barisan $\{ x_n \}_{n=1}^\infty$ yang konvergen ke $c$, atau dengan menerapkan
definisi kekontinuan secara langsung.  Rinciannya diserahkan kepada
pembaca.
\end{example}

\subsection{Komposisi fungsi kontinu}

Anda mungkin sudah menyadari bahwa salah satu alat dasar untuk
membangun fungsi yang rumit dari fungsi-fungsi sederhana adalah komposisi.
Ingat bahwa untuk dua fungsi $f$ dan $g$,
komposisi $f \circ g$ didefinisikan oleh
$(f \circ g)(x) \coloneqq f\bigl(g(x)\bigr)$.
Komposisi
fungsi-fungsi kontinu juga
kontinu.

\begin{prop} \label{prop:compositioncont}
Misalkan $A, B \subset \R$ serta $f \colon B \to \R$ dan $g \colon A \to B$ adalah
fungsi-fungsi.  Jika $g$ kontinu di $c \in A$ dan
$f$ kontinu di $g(c)$, maka $f \circ g \colon A \to \R$ kontinu
di $c$.
\end{prop}

\begin{proof}
Misalkan $\{ x_n \}_{n=1}^\infty$ suatu barisan dalam $A$ sedemikian sehingga
$\lim_{n\to\infty} x_n = c$.
Karena $g$ kontinu di $c$, barisan $\bigl\{ g(x_n) \bigr\}_{n=1}^\infty$ konvergen ke $g(c)$.
Karena $f$ kontinu di $g(c)$, barisan $\bigl\{ f\bigl(g(x_n)\bigr)
\bigr\}_{n=1}^\infty$ konvergen
ke $f\bigl(g(c)\bigr)$.
Jadi, $f \circ g$ kontinu di $c$.
\end{proof}

\begin{example}
Klaim: \emph{${\bigl(\sin(\nicefrac{1}{x})\bigr)}^2$ merupakan fungsi
kontinu pada $(0,\infty)$.}

Bukti: Fungsi $\nicefrac{1}{x}$ kontinu pada
$(0,\infty)$ dan $\sin(x)$ kontinu pada $(0,\infty)$ (bahkan
pada $\R$, tetapi $(0,\infty)$ merupakan daerah hasil $\nicefrac{1}{x}$).
Oleh karena itu, komposisi $\sin(\nicefrac{1}{x})$ kontinu.  Selain itu,
$x^2$ kontinu pada interval $[-1,1]$ (daerah hasil $\sin$).
Jadi, komposisi
${\bigl(\sin(\nicefrac{1}{x})\bigr)}^2$ kontinu pada $(0,\infty)$.
\end{example}

\subsection{Fungsi diskontinu}

Ketika $f$ tidak kontinu di $c$, kita
mengatakan bahwa $f$ \emph{\myindex{diskontinu}} di $c$, atau bahwa fungsi tersebut memiliki
\emph{\myindex{diskontinuitas}} di~$c$.
Proposisi berikut merupakan uji yang berguna dan langsung mengikuti
butir ketiga dari \propref{contbasic:prop}.

\begin{prop}
Misalkan $f \colon S \to \R$ suatu fungsi dan $c \in S$.  Andaikan
terdapat barisan $\{ x_n \}_{n=1}^\infty$ dengan $x_n \in S$ untuk semua $n$,
dan $\lim_{n\to\infty} x_n = c$
sedemikian sehingga $\bigl\{ f(x_n) \bigr\}_{n=1}^\infty$
tidak konvergen ke $f(c)$.
Maka $f$ diskontinu di $c$.
\end{prop}

Sekali lagi,
pernyataan bahwa
$\bigl\{ f(x_n) \bigr\}_{n=1}^\infty$ tidak konvergen ke $f(c)$
berarti bahwa barisan tersebut
entah sama sekali tidak konvergen,
atau konvergen ke sesuatu
yang berbeda dari $f(c)$.

\begin{example} \label{example:jumpdiscont}
Fungsi $f \colon \R \to \R$ yang didefinisikan oleh
\begin{equation*}
f(x) \coloneqq 
\begin{cases}
-1 & \text{jika } x < 0, \\
1 & \text{jika } x \geq 0
\end{cases}
\end{equation*}
tidak kontinu di 0.

Bukti: Tinjau $\{ \nicefrac{-1}{n} \}_{n=1}^\infty$, yang konvergen ke 0.  Maka
$f(\nicefrac{-1}{n}) = -1$ untuk setiap $n$,
sehingga
$\lim_{n\to\infty} f(\nicefrac{-1}{n}) = -1$, tetapi $f(0) = 1$.  Jadi, fungsi tersebut
tidak kontinu di 0.  Lihat
\figureref{fig:jumpdiscont}.

\begin{myfigureht}
\myincludegraphics{jumpdiscont}{%
Grafik fungsi yang bernilai minus 1 untuk x negatif
dan 1 untuk x nonnegatif.  Titik minus 1,
minus setengah, minus sepertiga, dan seterusnya ditandai pada
sumbu x dan nilai-nilai yang bersesuaian pada
grafik ditandai dengan titik-titik.}
\caption{Diskontinuitas loncat.  Nilai
$f(\nicefrac{-1}{n})$ dan $f(0)$ ditandai sebagai titik hitam.\label{fig:jumpdiscont}}
\end{myfigureht}

Perhatikan bahwa $f(\nicefrac{1}{n}) = 1$ untuk semua $n \in \N$.
Oleh karena itu, $\lim_{n\to\infty} f(\nicefrac{1}{n}) = f(0) = 1$.  Jadi,
$\bigl\{ f(x_n) \bigr\}_{n=1}^\infty$ dapat konvergen ke $f(0)$
untuk suatu barisan khusus
$\{ x_n \}_{n=1}^\infty$ yang menuju 0,
meskipun fungsi tersebut diskontinu di 0.

Terakhir, tinjau $f\Bigl(\frac{{(-1)}^n}{n}\Bigr) = {(-1)}^n$.
Barisan ini divergen.
\end{example}

\begin{example}
Sebagai contoh ekstrem, tinjau
\emph{\myindex{fungsi Dirichlet}}\footnote{Dinamai menurut matematikawan
Jerman
\href{https://en.wikipedia.org/wiki/Peter_Gustav_Lejeune_Dirichlet}{Johann Peter Gustav Lejeune
Dirichlet}
(1805--1859).}.
\begin{equation*}
f(x) \coloneqq
\begin{cases}
1 & \text{jika } x \text{ rasional,} \\
0 & \text{jika } x \text{ irasional.}
\end{cases}
\avoidbreak
\end{equation*}
Fungsi $f$ diskontinu di setiap $c \in \R$.

Bukti:
Jika $c$ rasional, ambil barisan $\{ x_n \}_{n=1}^\infty$
bilangan irasional sedemikian sehingga $\lim_{n\to\infty} x_n = c$ (mengapa hal ini mungkin?).  Maka $f(x_n) = 0$
sehingga $\lim_{n\to\infty} f(x_n) = 0$, tetapi $f(c) = 1$.
Jika $c$ irasional, ambil barisan
$\{ x_n \}_{n=1}^\infty$
bilangan rasional
yang konvergen ke $c$ (mengapa hal ini mungkin?).  Maka $\lim_{n\to\infty} f(x_n) = 1$, tetapi $f(c) = 0$.
\end{example}

Mari kita uji batas intuisi kita.  Mungkinkah
terdapat fungsi yang kontinu di semua bilangan irasional, tetapi
diskontinu di semua bilangan rasional?  Di dekat setiap bilangan irasional terdapat bilangan rasional
dengan jarak sekecil apa pun.  Mungkin agak mengejutkan,
jawabannya adalah ya, fungsi demikian ada.  Contoh berikut disebut
\emph{\myindex{fungsi Thomae}}\footnote{Dinamai menurut matematikawan
Jerman
\href{https://en.wikipedia.org/wiki/Carl_Johannes_Thomae}{Carl Johannes Thomae}
(1840--1921).} atau
\emph{\myindex{fungsi popcorn}}.

\begin{example} \label{popcornfunction:example}
Definisikan $f \colon (0,1) \to \R$ dengan
\begin{equation*}
f(x) \coloneqq 
\begin{cases}
\nicefrac{1}{k} & \begin{aligned}
&\text{jika } x=\nicefrac{m}{k},\ m,k \in \N, \\ &\text{dan } m,k \text{ relatif prima (bentuk paling sederhana),}\end{aligned} \\
0 & \text{jika } x \text{ irasional.}
\end{cases}
\end{equation*}
Lihat grafik $f$
dalam \figureref{popcornfig}.
Kita akan menunjukkan bahwa
$f$ kontinu di setiap $c$ irasional dan
diskontinu di setiap $c$ rasional.
\begin{myfigureht}
\myincludegraphics{popcornfig}{%
Grafik suatu fungsi yang tersusun dari titik-titik.
Titik-titik ini tampak tersusun membentuk
berbagai segitiga, tetapi jelas bahwa
fungsi ini diskontinu di semua titik tersebut,
dan terdapat semakin banyak titik semacam itu
ketika kita semakin dekat ke sumbu x.}
\caption{Grafik \myquote{fungsi popcorn.}\label{popcornfig}}
\end{myfigureht}

Bukti:
Misalkan $c = \nicefrac{m}{k}$ rasional dan dalam bentuk paling sederhana.  Ambil barisan
bilangan irasional $\{ x_n \}_{n=1}^\infty$ sedemikian sehingga $\lim_{n\to\infty} x_n = c$.  Maka
$\lim_{n\to\infty} f(x_n) = \lim_{n\to\infty} 0 = 0$, tetapi $f(c) = \nicefrac{1}{k} \neq 0$.  Jadi, $f$
diskontinu di $c$.

Sekarang misalkan $c$ irasional, sehingga $f(c) = 0$.  Ambil barisan
$\{ x_n \}_{n=1}^\infty$ dalam $(0,1)$ sedemikian sehingga $\lim_{n\to\infty} x_n = c$.
Untuk $\epsilon > 0$, pilih $K \in \N$ sedemikian
sehingga $\nicefrac{1}{K} < \epsilon$
berdasarkan \hyperref[thm:arch:i]{sifat Archimedes}.
Jika $\nicefrac{m}{k} \in (0,1)$ dan $m,k \in \N$, maka $0 < m < k$.
Jadi, hanya terdapat berhingga banyak bilangan rasional dalam $(0,1)$
yang penyebut $k$-nya dalam bentuk paling sederhana kurang dari $K$.
Karena $\lim_{n\to\infty} x_n = c$, setiap bilangan yang tidak sama dengan $c$ hanya dapat muncul
berhingga kali dalam $\{ x_n \}_{n=1}^\infty$.
Oleh karena itu,
terdapat $M$ sedemikian sehingga untuk $n \geq M$, semua bilangan $x_n$
yang rasional
memiliki penyebut yang lebih besar daripada atau sama dengan $K$.  Jadi, untuk $n \geq M$,
\begin{equation*}
\babs{f(x_n) - 0} = f(x_n) \leq \nicefrac{1}{K} < \epsilon .
\end{equation*}
Oleh karena itu, $f$ kontinu di titik irasional $c$.
\end{example}

Mari kita akhiri dengan contoh yang lebih mudah.

\begin{example} \label{example:removablediscont}
Definisikan
$g \colon \R \to \R$ dengan $g(x) \coloneqq 0$ jika $x \neq 0$ dan
$g(0) \coloneqq 1$.  Maka $g$ tidak kontinu di nol, tetapi kontinu di semua titik lain (mengapa?).
Titik $x=0$ disebut \emph{\myindex{diskontinuitas yang dapat dihilangkan}}.  Hal itu
karena jika kita mengubah definisi $g$ dengan menetapkan
$g(0)$ sama dengan $0$, kita akan memperoleh fungsi kontinu.  Di sisi lain,
misalkan $f$ adalah fungsi pada \exampleref{example:jumpdiscont}.
Maka $f$ tidak memiliki
diskontinuitas yang dapat dihilangkan di $0$.  Bagaimanapun kita mendefinisikan $f(0)$, fungsi tersebut
tetap tidak kontinu.  Perbedaannya ialah
$\lim_{x\to 0} g(x)$ ada, sedangkan
$\lim_{x\to 0} f(x)$ tidak ada.

Kita tetap menggunakan contoh ini untuk menunjukkan fenomena lain.
Misalkan $A \coloneqq \{ 0 \}$, maka $g|_A$ kontinu (mengapa?), sedangkan $g$ tidak kontinu pada $A$.
Demikian pula, jika $B \coloneqq \R \setminus \{0 \}$, maka $g|_B$ juga kontinu,
dan sebenarnya $g$ kontinu pada $B$.
\end{example}

\subsection{Latihan}

\begin{exercise}
Dengan menggunakan definisi kekontinuan secara langsung, buktikan bahwa
$f \colon \R \to \R$ yang didefinisikan oleh
$f(x) \coloneqq x^2$ bersifat kontinu.
\end{exercise}

\begin{exercise}
Dengan menggunakan definisi kekontinuan secara langsung, buktikan bahwa
$f \colon (0,\infty) \to \R$ yang didefinisikan oleh
$f(x) \coloneqq \nicefrac{1}{x}$ bersifat kontinu.
\end{exercise}

\begin{exercise}
Definisikan $f \colon \R \to \R$ dengan
\begin{equation*}
f(x) \coloneqq
\begin{cases}
x & \text{jika } x \text{ rasional,} \\
x^2 & \text{jika } x \text{ irasional.}
\end{cases}
\avoidbreak
\end{equation*}
Dengan menggunakan definisi kekontinuan secara langsung, buktikan bahwa
$f$ kontinu di $1$ dan diskontinu di $2$.
\end{exercise}

\begin{exercise}
Definisikan $f \colon \R \to \R$ dengan
\begin{equation*}
f(x) \coloneqq
\begin{cases}
\sin(\nicefrac{1}{x}) & \text{jika } x \neq 0, \\
0 & \text{jika } x=0.
\end{cases}
\avoidbreak
\end{equation*}
Apakah $f$ kontinu?  Buktikan jawaban Anda.
\end{exercise}

\begin{exercise}
Definisikan $f \colon \R \to \R$ dengan
\begin{equation*}
f(x) \coloneqq
\begin{cases}
x \sin(\nicefrac{1}{x}) & \text{jika } x \neq 0, \\
0 & \text{jika } x=0.
\end{cases}
\avoidbreak
\end{equation*}
Apakah $f$ kontinu?  Buktikan jawaban Anda.
\end{exercise}

\begin{exercise}
Buktikan \propref{contalg:prop}.
\end{exercise}

\begin{exercise} \label{exercise:restrictioncontinuous}
Misalkan $S \subset \R$ dan $A \subset S$.  Misalkan $f \colon S \to \R$
suatu fungsi kontinu.
Buktikan bahwa pembatasan $f|_A$ bersifat kontinu.
\end{exercise}

\begin{exercise}
Andaikan $S \subset \R$ sedemikian sehingga $(c-\alpha,c+\alpha) \subset S$ untuk suatu $c \in \R$
dan $\alpha > 0$.
Misalkan $f \colon S \to \R$ suatu fungsi dan $A \coloneqq (c-\alpha,c+\alpha)$.  Buktikan bahwa
jika $f|_A$ kontinu di $c$, maka $f$ kontinu di $c$.
\end{exercise}

\begin{exercise}
Berikan contoh fungsi $f \colon \R \to \R$ dan $g \colon \R \to \R$
sedemikian sehingga fungsi $h$, yang didefinisikan oleh $h(x) \coloneqq f(x) + g(x)$, bersifat kontinu,
tetapi $f$ dan $g$ tidak kontinu.  Dapatkah Anda menemukan $f$ dan $g$ yang tidak kontinu
di titik mana pun, tetapi $h$ merupakan fungsi kontinu?
\end{exercise}

\begin{exercise}
Misalkan $f \colon \R \to \R$ dan 
$g \colon \R \to \R$ fungsi-fungsi kontinu.  Andaikan bahwa
$f(r) = g(r)$ berlaku untuk semua $r \in \Q$.
Tunjukkan bahwa $f(x) = g(x)$ untuk semua $x \in \R$.
\end{exercise}

\begin{exercise} \label{exercise:positivecontneigh}
Misalkan $f \colon \R \to \R$ kontinu.  Andaikan $f(c) > 0$.  Tunjukkan bahwa
terdapat $\alpha > 0$ sedemikian sehingga untuk semua $x \in (c-\alpha,c+\alpha)$,
berlaku $f(x) > 0$.
\end{exercise}

\begin{exercise}
Misalkan $f \colon \Z \to \R$ suatu fungsi.  Tunjukkan bahwa $f$ kontinu.
\end{exercise}

\begin{exercise} \label{exercise:contseqalt}
Misalkan $f \colon S \to \R$ suatu fungsi dan $c \in S$ sedemikian sehingga untuk setiap
barisan $\{ x_n \}_{n=1}^\infty$ dalam $S$ dengan $\lim_{n\to\infty} x_n = c$, barisan
$\bigl\{ f(x_n) \bigr\}_{n=1}^\infty$ konvergen.  Tunjukkan bahwa $f$ kontinu di $c$.
\end{exercise}

\begin{exercise}
Andaikan $f \colon [-1,0] \to \R$ dan $g \colon [0,1] \to \R$ kontinu
serta $f(0) = g(0)$.  Definisikan $h \colon [-1,1] \to \R$ dengan 
$h(x) \coloneqq f(x)$ jika $x \leq 0$ dan $h(x) \coloneqq g(x)$ jika $x > 0$.  Tunjukkan bahwa
$h$ kontinu.
\end{exercise}

\begin{exercise}
Andaikan $g \colon \R \to \R$ suatu fungsi kontinu sedemikian sehingga $g(0) = 0$,
dan andaikan $f \colon \R \to \R$ memenuhi
$\babs{f(x)-f(y)} \leq g(x-y)$ untuk semua $x$ dan $y$.  Tunjukkan bahwa $f$
kontinu.
\end{exercise}

\begin{exercise}[Menantang]
Andaikan $f \colon \R \to \R$ kontinu di $0$ dan
memenuhi $f(x+y) = f(x) + f(y)$ untuk setiap $x$ dan $y$.
Tunjukkan bahwa $f(x) = ax$ untuk suatu $a \in \R$.
Petunjuk: Tunjukkan bahwa $f(nx) = nf(x)$, lalu tunjukkan bahwa $f$ kontinu pada $\R$.
Kemudian tunjukkan bahwa $\nicefrac{f(x)}{x} = f(1)$ untuk setiap bilangan rasional $x \neq 0$.
\end{exercise}

\begin{exercise} \label{exercise:minmaxcont}
Andaikan $S \subset \R$ dan
misalkan $f \colon S \to \R$ dan
$g \colon S \to \R$ fungsi-fungsi kontinu.
Definisikan $p \colon S \to \R$ dengan
$p(x) \coloneqq \max \bigl\{ f(x) , g(x) \bigr\}$ dan
$q \colon S \to \R$ dengan
$q(x) \coloneqq \min \bigl\{ f(x) , g(x) \bigr\}$.  Buktikan bahwa $p$ dan $q$
kontinu.
\end{exercise}

\begin{exercise}
Andaikan $f \colon [-1,1] \to \R$ suatu fungsi yang kontinu di semua $x \in
[-1,1] \setminus \{ 0 \}$.  Tunjukkan bahwa untuk setiap $\epsilon$ sedemikian
sehingga $0 < \epsilon < 1$, terdapat
fungsi $g \colon [-1,1] \to \R$ yang kontinu pada seluruh $[-1,1]$ sedemikian sehingga
$f(x) = g(x)$ untuk semua $x \in [-1,-\epsilon] \cup [\epsilon,1]$, serta 
$\babs{g(x)} \leq \babs{f(x)}$ untuk semua $x \in [-1,1]$.
\end{exercise}

\begin{exercise}[Menantang]
Suatu fungsi $f \colon I \to \R$ disebut \emph{\myindex{konveks}} jika
setiap kali $a \leq x \leq b$ untuk $a,x,b$ dalam $I$ dengan $a<b$, berlaku
$f(x) \leq f(a) \frac{b-x}{b-a} + f(b) \frac{x-a}{b-a}$.  Dengan kata lain,
garis yang ditarik antara $\bigl(a,f(a)\bigr)$ dan $\bigl(b,f(b)\bigr)$ 
berada di atas grafik $f$.
\begin{enumerate}[a)]
\item
Buktikan bahwa
jika $I = (\alpha,\beta)$ merupakan interval terbuka dan $f \colon I \to \R$ konveks,
maka $f$ kontinu.
\item
Carilah contoh fungsi konveks $f \colon [0,1] \to \R$ yang
tidak kontinu.
\end{enumerate}
\end{exercise}


%%%%%%%%%%%%%%%%%%%%%%%%%%%%%%%%%%%%%%%%%%%%%%%%%%%%%%%%%%%%%%%%%%%%%%%%%%%%%%

\sectionnewpage
\section{Teorema nilai ekstrem dan nilai antara}
\label{sec:minmaxint}

%mbxINTROSUBSECTION

\sectionnotes{1,5 kuliah}

Fungsi kontinu pada interval tertutup dan terbatas
berperilaku sangat baik.

\subsection{Teorema min--maks atau teorema nilai ekstrem}

Ingat bahwa $f \colon [a,b] \to \R$ disebut
\emph{terbatas}\index{fungsi terbatas}\index{fungsi!terbatas}
jika terdapat $B \in \R$ sedemikian sehingga
$\babs{f(x)} \leq B$ untuk semua $x \in [a,b]$.
Untuk fungsi kontinu pada interval tertutup dan terbatas, kita mempunyai lemma berikut.

\begin{lemma}
Fungsi kontinu $f \colon [a,b] \to \R$ bersifat terbatas.
\end{lemma}

\begin{proof}
Kita membuktikan pernyataan tersebut dengan kontraposisi.  Misalkan $f$ tidak terbatas.
Maka untuk setiap
$n \in \N$, terdapat $x_n \in [a,b]$ sedemikian sehingga
\begin{equation*}
\babs{f(x_n)} \geq n .
\end{equation*}
Barisan $\{ x_n \}_{n=1}^\infty$ terbatas karena $a \leq x_n \leq b$.
Berdasarkan \hyperref[thm:bwseq]{teorema Bolzano--Weierstrass},
terdapat subbarisan konvergen $\{ x_{n_i} \}_{i=1}^\infty$.
Misalkan $x \coloneqq \lim_{i\to\infty} x_{n_i}$.
Karena $a \leq x_{n_i} \leq b$ untuk semua $i$, berlaku $a \leq x \leq b$.
Barisan $\bigl\{ f(x_{n_i}) \bigr\}_{i=1}^\infty$ tidak terbatas
karena
$\babs{f(x_{n_i})} \geq n_i \geq i$.
Dengan demikian, $f$ tidak kontinu di $x$ karena
\begin{equation*}
f(x)
=
f\Bigl( \lim_{i\to\infty} x_{n_i} \Bigr) ,
\qquad \text{tetapi} \qquad
\lim_{i\to\infty} f(x_{n_i}) \enspace \text{tidak ada.} \qedhere
\end{equation*}
\end{proof}

Perhatikan satu hal penting dalam pembuktian tersebut.
Keterbatasan $[a,b]$ memungkinkan kita menggunakan Bolzano--Weierstrass,
sedangkan sifat tertutupnya memastikan bahwa limit kembali berada dalam $[a,b]$.
Teknik ini lazim digunakan: Cari barisan dengan sifat tertentu,
lalu gunakan Bolzano--Weierstrass untuk memperoleh subbarisan yang juga konvergen.

Ingat dari kalkulus bahwa $f \colon S \to \R$ mencapai
\emph{\myindex{minimum mutlak}}\index{minimum!mutlak}%
\index{mencapai minimum mutlak}
di $c \in S$ jika
\begin{equation*}
f(x) \geq f(c) \qquad \text{untuk semua } x \in S.
\end{equation*}
Di sisi lain, $f$ mencapai
\emph{\myindex{maksimum mutlak}}\index{maksimum!mutlak}%
\index{mencapai maksimum mutlak}
di $c \in S$ jika
\begin{equation*}
f(x) \leq f(c) \qquad \text{untuk semua } x \in S.
\end{equation*}
Jika terdapat $c \in S$ semacam itu, kita mengatakan
$f$ \emph{mencapai minimum mutlak (atau\ maksimum mutlak) pada
$S$}, dan kita menyebut
$f(c)$ sebagai \emph{minimum mutlak (atau\ maksimum mutlak)}.
%See \figureref{fig:minmax}.
\begin{myfigureht}
\myincludegraphics{minmax}{%
Grafik suatu fungsi pada interval dari a hingga b.
Grafik tersebut berada di antara dua garis horizontal,
satu ditandai sebagai maksimum mutlak f yang sama dengan
f dari c, dan
satu ditandai sebagai minimum mutlak f yang sama dengan
f dari d.}
\caption{$f \colon [a,b] \to \R$ mencapai maksimum mutlak $f(c)$ di
$c$, dan minimum mutlak $f(d)$ di $d$.\label{fig:minmax}}
\end{myfigureht}

Jika $S$ merupakan interval tertutup
dan terbatas, maka fungsi kontinu $f$
tidak hanya terbatas, tetapi juga harus mencapai minimum mutlak dan maksimum
mutlak pada $S$.

\begin{thm}[Teorema minimum--maksimum / Teorema nilai ekstrem]
\index{teorema minimum-maksimum}%
\index{teorema maksimum-minimum}%
\index{teorema nilai ekstrem}%
Fungsi kontinu $f \colon [a,b] \to \R$
mencapai minimum mutlak dan maksimum mutlak pada $[a,b]$.
\end{thm}

Sekali lagi, penting diperhatikan bahwa domain $f$ adalah interval tertutup dan terbatas $[a,b]$.

\begin{proof}
Lemma tersebut menyatakan bahwa $f$ terbatas, sehingga
himpunan $f\bigl([a,b]\bigr) = \bigl\{ f(x) : x \in [a,b] \bigr\}$ memiliki supremum dan infimum.
Terdapat barisan-barisan
dalam himpunan $f\bigl([a,b]\bigr)$ yang mendekati supremum dan infimumnya.
Artinya, terdapat barisan
$\bigl\{ f(x_n) \bigr\}_{n=1}^\infty$ dan $\bigl\{ f(y_n)
\bigr\}_{n=1}^\infty$, dengan $x_n$ dan $y_n$ berada dalam $[a,b]$,
sedemikian sehingga
\begin{equation*}
\lim_{n\to\infty} f(x_n) = \inf f\bigl([a,b]\bigr) \qquad \text{dan} \qquad
\lim_{n\to\infty} f(y_n) = \sup f\bigl([a,b]\bigr).
\end{equation*}
Kita belum selesai; kita perlu menentukan tempat minimum dan maksimum tersebut dicapai.
Masalahnya, barisan $\{ x_n \}_{n=1}^\infty$ dan
$\{ y_n \}_{n=1}^\infty$ belum tentu konvergen.
Kita mengetahui bahwa $\{ x_n \}_{n=1}^\infty$ dan $\{ y_n \}_{n=1}^\infty$ terbatas
(unsur-unsurnya berada dalam interval terbatas $[a,b]$).
Terapkan
\hyperref[thm:bwseq]{teorema Bolzano--Weierstrass}
untuk memperoleh
subbarisan konvergen
$\{ x_{n_i} \}_{i=1}^\infty$ dan
$\{ y_{m_i} \}_{i=1}^\infty$.  Misalkan
\begin{equation*}
x \coloneqq \lim_{i\to\infty} x_{n_i}
\qquad \text{dan} \qquad
y \coloneqq \lim_{i\to\infty} y_{m_i}.
\end{equation*}
Karena $a \leq x_{n_i} \leq b$ untuk semua $i$, berlaku $a \leq x \leq b$.
Demikian pula, $a \leq y \leq b$.  Jadi, $x$ dan $y$ berada dalam $[a,b]$.
Untuk barisan yang konvergen,
limit suatu subbarisan sama dengan limit
barisan tersebut.  Selain itu, kita dapat melewatkan limit melalui fungsi kontinu $f$.  Jadi,
\begin{equation*}
\inf f\bigl([a,b]\bigr) = \lim_{n\to\infty} f(x_n)
= \lim_{i\to\infty} f(x_{n_i}) = 
f \Bigl( \lim_{i\to\infty} x_{n_i} \Bigr) = f(x) .
\end{equation*}
Demikian pula,
\begin{equation*}
\sup f\bigl([a,b]\bigr) = \lim_{n\to\infty} f(y_n)
= \lim_{i\to\infty} f(y_{m_i}) = 
f \Bigl( \lim_{i\to\infty} y_{m_i} \Bigr) = f(y) .
\end{equation*}
Dengan demikian, $f$ mencapai minimum mutlak di $x$ dan
maksimum mutlak di $y$.
\end{proof}

\begin{example}
Fungsi $f(x) \coloneqq x^2+1$ yang didefinisikan pada interval $[-1,2]$ mencapai minimum
di $x=0$ dengan $f(0) = 1$.  Fungsi tersebut mencapai maksimum di $x=2$ dengan $f(2) = 5$.
Perhatikan bahwa domain definisi berpengaruh.  Jika kita memilih domain
$[-10,10]$, maka $f$ tidak lagi mencapai maksimum di $x=2$.
Sebaliknya,
maksimum dicapai di $x=10$ dan juga di $x=-10$.
\end{example}

Melalui contoh-contoh, kita menunjukkan bahwa berbagai hipotesis teorema tersebut
benar-benar diperlukan.

\begin{example}
Fungsi $f \colon \R \to \R$ yang didefinisikan oleh $f(x) \coloneqq x$
tidak mencapai minimum maupun maksimum.  Jadi, penting bahwa
kita meninjau interval terbatas.
\end{example}

\begin{example}
Fungsi $f \colon (0,1) \to \R$ yang didefinisikan oleh $f(x) \coloneqq \nicefrac{1}{x}$
tidak mencapai minimum maupun maksimum.
Fungsi tersebut kontinu, tetapi $(0,1)$ tidak tertutup.
Nilai-nilai fungsi tersebut
tak terbatas ketika kita mendekati $0$.  Selain itu, ketika kita mendekati $x=1$, nilai-nilai
fungsi mendekati $1$, tetapi $f(x) > 1$ untuk semua $x \in (0,1)$.  Tidak ada
$x \in (0,1)$ sedemikian sehingga $f(x) = 1$.  Jadi, penting bahwa
kita meninjau interval tertutup.
\end{example}

\begin{example}
Kekontinuan juga penting.
Definisikan $f \colon [0,1] \to \R$ dengan
$f(x) \coloneqq \nicefrac{1}{x}$ untuk $x > 0$ dan tetapkan $f(0) \coloneqq 0$.
Fungsi tersebut tidak mencapai maksimum.  Domain $[0,1]$ tertutup dan
terbatas, tetapi masalahnya adalah
fungsi tersebut tidak kontinu di 0.
\end{example}

\subsection{Teorema nilai antara Bolzano}

Teorema nilai antara Bolzano merupakan salah satu landasan utama analisis.
Teorema ini kadang-kadang hanya disebut teorema nilai antara atau cukup
teorema Bolzano.  Untuk membuktikan teorema Bolzano, kita membuktikan
lema yang lebih sederhana berikut ini.

\begin{lemma} \label{IVT:lemma}
Misalkan $f \colon [a,b] \to \R$ suatu fungsi kontinu.
Andaikan $f(a) < 0$ dan $f(b) > 0$.
Maka terdapat bilangan $c \in (a,b)$
sedemikian sehingga $f(c) = 0$.
\end{lemma}

\begin{proof}
Kita mendefinisikan dua barisan $\{ a_n \}_{n=1}^\infty$
dan $\{ b_n \}_{n=1}^\infty$ secara induktif:
\begin{enumerate}[(i)]
\item Tetapkan $a_1 \coloneqq a$ dan $b_1 \coloneqq b$.
\item Jika $f\left(\frac{a_n+b_n}{2}\right) \geq 0$, tetapkan $a_{n+1} \coloneqq a_n$ dan
$b_{n+1} \coloneqq \frac{a_n+b_n}{2}$.
\item Jika $f\left(\frac{a_n+b_n}{2}\right) < 0$, tetapkan $a_{n+1} \coloneqq \frac{a_n+b_n}{2}$ dan
$b_{n+1} \coloneqq b_n$.
\end{enumerate}
\begin{myfigureht}
\myincludegraphics{bisect}{%
Grafik suatu fungsi yang memotong
sumbu-x di c sambil naik.  Interval
a subskrip 1 hingga b subskrip 1 ditandai; f bernilai negatif
di a subskrip 1 dan positif di b subskrip 1.
Interval berikutnya berawal di a subskrip 2, yang sama dengan a subskrip 1,
dan b subskrip 2 berada di tengah interval sebelumnya.
Sekali lagi, fungsi bernilai negatif di a subskrip 2 dan positif
di b subskrip 2.  Kita melanjutkan dengan interval-interval
tempat f bernilai negatif di sebelah kiri dan positif
di sebelah kanan.
Interval berikutnya berawal di a subskrip 3, yang sama dengan
a subskrip 2, dan berakhir di b subskrip 3, yang
berada di tengah interval sebelumnya.
Untuk interval berikutnya, a subskrip 4 berada di tengah dan
b subskrip 4 sama dengan b subskrip 3.
Berikutnya, a subskrip 5 kembali berada di tengah dan b subskrip 5
sama dengan b subskrip 4.  Bilangan c berada
di dalam semua interval ini.}
\caption{Menentukan akar (metode biseksi).\label{bisectfig}}
\end{myfigureht}
Lihat \figureref{bisectfig} untuk contoh lima langkah pertama.
Jika $a_n < b_n$, maka $a_n < \frac{a_n+b_n}{2} < b_n$.  Jadi,
$a_{n+1} < b_{n+1}$.
Karena $a_1 = a < b = b_1$,
\hyperref[induction:thm]{induksi} memberikan
$a_n < b_n$ untuk semua $n$.
Selain itu, $a_n \leq a_{n+1}$ dan
$b_n \geq b_{n+1}$ untuk semua $n$; dengan kata lain, kedua barisan tersebut monoton.
Karena $a_n < b_n \leq b_1 = b$ dan
$b_n > a_n \geq a_1 = a$ untuk semua $n$,
kedua barisan tersebut juga terbatas.  Oleh karena itu,
kedua barisan tersebut konvergen.
Misalkan $c \coloneqq \lim_{n\to\infty} a_n$ dan $d \coloneqq \lim_{n\to\infty} b_n$;
berlaku pula $a \leq c \leq d \leq b$.  Kita perlu
menunjukkan bahwa $c=d$.
Perhatikan bahwa
\begin{equation*}
b_{n+1} - a_{n+1} = \frac{b_n-a_n}{2}.
\end{equation*}
Dengan \hyperref[induction:thm]{induksi},
\begin{equation*}
b_n - a_n = \frac{b_1-a_1}{2^{n-1}} = 2^{1-n} (b-a) .
\end{equation*}
Karena $2^{1-n}(b-a)$ konvergen ke nol, dengan mengambil limit ketika $n$ menuju
tak hingga kita memperoleh
\begin{equation*}
d-c = \lim_{n\to\infty} (b_n - a_n) =
\lim_{n\to\infty} 2^{1-n} (b-a) = 0.
\avoidbreak
\end{equation*}
Dengan kata lain, $c=d$.

Berdasarkan konstruksinya, untuk semua $n$,
\begin{equation*}
f(a_n) < 0
\qquad \text{dan} \qquad
f(b_n) \geq 0 .
\end{equation*}
Karena
$\lim_{n\to\infty} a_n = \lim_{n\to\infty} b_n = c$
dan $f$ kontinu di $c$, kita dapat mengambil
limit dalam ketaksamaan-ketaksamaan tersebut:
\begin{equation*}
f(c) = \lim_{n\to\infty} f(a_n) \leq 0
\qquad \text{dan} \qquad
f(c) = \lim_{n\to\infty} f(b_n) \geq 0 .
\end{equation*}
Karena $f(c) \geq 0$ dan
$f(c) \leq 0$, kita menyimpulkan bahwa $f(c) = 0$.
Dengan demikian, $c \neq a$ dan $c \neq b$, sehingga
$a < c < b$.
\end{proof}

\begin{thm}[Teorema nilai antara Bolzano] \label{IVT:thm}
\index{teorema Bolzano}
\index{teorema nilai antara Bolzano}
\index{teorema nilai antara}
Misalkan $f \colon [a,b] \to \R$ suatu fungsi kontinu.
Andaikan $y \in \R$ sedemikian sehingga $f(a) < y < f(b)$
atau $f(a) > y > f(b)$.  Maka terdapat $c \in (a,b)$
sedemikian sehingga $f(c) = y$.
\end{thm}

Teorema tersebut menyatakan bahwa suatu fungsi kontinu pada interval tertutup
mencapai semua nilai di antara nilai-nilai pada kedua titik ujungnya.

\begin{proof}
Jika $f(a) < y < f(b)$, definisikan $g(x) \coloneqq f(x)-y$.  Dengan demikian,
$g(a) < 0$ dan $g(b) > 0$, lalu kita menerapkan \lemmaref{IVT:lemma}
pada $g$ untuk memperoleh $c$.  Jika $g(c) = 0$, maka $f(c) = y$.

Demikian pula, jika $f(a) > y > f(b)$, definisikan $g(x) \coloneqq y-f(x)$.
Sekali lagi, $g(a) < 0$ dan $g(b) > 0$, lalu kita menerapkan \lemmaref{IVT:lemma} untuk
memperoleh $c$.
Seperti sebelumnya, jika $g(c) = 0$, maka $f(c) = y$.
\end{proof}

Jika suatu fungsi kontinu, pembatasannya
pada suatu himpunan bagian juga kontinu; jika $f \colon S \to \R$ kontinu dan
$[a,b] \subset S$, maka $f|_{[a,b]}$ juga kontinu.  Pada umumnya, kita
menerapkan teorema tersebut pada suatu fungsi yang kontinu pada himpunan $S$ yang lebih besar,
tetapi membatasi perhatian kita pada suatu interval.


Pembuktian lema tersebut memberi tahu kita cara menentukan akar $c$.
Pembuktian itu tidak hanya berguna bagi kita para matematikawan murni,
tetapi juga memuat gagasan yang berguna dalam matematika terapan,
yang disebut \emph{\myindex{metode biseksi}}.


\begin{example}[Metode biseksi] %\index{metode biseksi}
Polinom $f(x) \coloneqq x^3-2x^2+x-1$ memiliki akar real dalam $(1,2)$.  Kita cukup
memperhatikan bahwa $f(1) = -1$ dan $f(2) = 1$.  Oleh karena itu, harus terdapat titik $c
\in (1,2)$ sedemikian sehingga $f(c) = 0$.  Untuk memperoleh hampiran akar yang
lebih baik, kita mengikuti pembuktian \lemmaref{IVT:lemma}.
Kita memeriksa 1.5 dan memperoleh $f(1.5) = -0.625$.  Oleh karena itu,
terdapat akar polinom tersebut dalam $(1.5,2)$.  Selanjutnya, kita memeriksa 1.75
dan mendapati bahwa $f(1.75) \approx -0.016$.  Jadi, terdapat akar $f$ dalam
$(1.75,2)$.  Selanjutnya, kita memeriksa 1.875 dan memperoleh $f(1.875) \approx 0.44$;
dengan demikian, terdapat akar dalam $(1.75,1.875)$.  Kita mengikuti prosedur ini hingga memperoleh
ketelitian yang memadai.  Bahkan, akar tersebut berada di $c \approx 1.7549$.
\end{example}

Teknik tersebut merupakan metode paling sederhana untuk menentukan akar polinom,
suatu masalah yang umum dalam matematika terapan.
Secara umum, menentukan akar secara cepat, tepat,
dan otomatis merupakan hal yang sulit.
Ada metode lain yang lebih cepat untuk menentukan akar polinom, seperti
metode Newton.  Salah satu keunggulan metode di atas adalah
kesederhanaannya.  Begitu kita menemukan suatu interval tempat teorema nilai antara
berlaku, kita dijamin dapat menemukan hampiran akar dengan ketelitian yang
diinginkan dalam sejumlah langkah berhingga.  Selain itu, metode
biseksi menentukan akar fungsi kontinu apa pun,
bukan hanya polinom.

Teorema tersebut menjamin adanya suatu $c$ sedemikian sehingga $f(c) = y$, tetapi
mungkin terdapat akar lain dari persamaan $f(c) = y$.  Jika kita mengikuti
prosedur dalam pembuktian, kita dijamin dapat menemukan hampiran untuk
salah satu akar tersebut.  Kita perlu berusaha lebih keras untuk menemukan akar-akar lainnya.

\medskip

Polinom berderajat genap mungkin tidak memiliki akar real sama sekali.
Tidak ada bilangan real $x$ sedemikian sehingga $x^2+1 = 0$.  Sebaliknya,
polinom berderajat ganjil selalu memiliki setidaknya satu akar real.

\begin{prop}
Misalkan $f(x)$ suatu polinom berderajat ganjil.  Maka $f$ memiliki akar real.
\end{prop}

\begin{proof}
Andaikan $f$ suatu polinom berderajat ganjil $d$.  Kita tuliskan
\begin{equation*}
f(x) = a_d x^d + a_{d-1} x^{d-1} + \cdots + a_1 x + a_0 ,
\end{equation*}
dengan $a_d \neq 0$.  Kita membaginya dengan $a_d$ untuk memperoleh
\emph{\myindex{polinom monik}}\footnote{Istilah \emph{monik} berarti bahwa
koefisien $x^d$ adalah 1.}
\begin{equation*}
g(x) \coloneqq x^d + b_{d-1} x^{d-1} + \cdots + b_1 x + b_0 ,
\end{equation*}
dengan $b_k = \nicefrac{a_k}{a_d}$.
Mari kita tunjukkan bahwa $g(n)$
positif untuk suatu $n \in \N$ yang cukup besar.
Pertama, kita membandingkan suku berderajat tertinggi dengan suku-suku lainnya:
\begin{equation*}
\begin{split}
\abs{\frac{b_{d-1} n^{d-1} + \cdots + b_1 n + b_0}{n^d}}
& =
\frac{\sabs{b_{d-1} n^{d-1} + \cdots + b_1 n + b_0}}{n^d}
\\
& \leq
\frac{\sabs{b_{d-1}} n^{d-1} + \cdots + \sabs{b_1} n + \sabs{b_0}}{n^d}
\\
& \leq
\frac{\sabs{b_{d-1}} n^{d-1} + \cdots + \sabs{b_1} n^{d-1} + \sabs{b_0} n^{d-1}}{n^d}
\\
& =
\frac{n^{d-1}\bigl(\sabs{b_{d-1}} + \cdots + \sabs{b_1} + \sabs{b_0}\bigr)}{n^d}
\\
& =
\frac{1}{n}
\bigl(\sabs{b_{d-1}} + \cdots + \sabs{b_1} + \sabs{b_0}\bigr) .
\end{split}
\end{equation*}
Oleh karena itu,
\begin{equation*}
\lim_{n\to\infty} \frac{b_{d-1} n^{d-1} + \cdots + b_1 n + b_0}{n^d}
= 0 .
\end{equation*}
Jadi, terdapat $M \in \N$ sedemikian sehingga
\begin{equation*}
\abs{\frac{b_{d-1} M^{d-1} + \cdots + b_1 M + b_0}{M^d}} < 1 ,
\end{equation*}
yang mengakibatkan
\begin{equation*}
-(b_{d-1} M^{d-1} + \cdots + b_1 M + b_0) < M^d .
\end{equation*}
Oleh karena itu, $g(M) > 0$.

Selanjutnya, tinjau $g(-n)$ untuk $n \in \N$.  Dengan argumen serupa,
terdapat $K \in \N$ sedemikian sehingga
$b_{d-1} {(-K)}^{d-1} + \cdots + b_1 (-K) + b_0 < K^d$
dan karena itu $g(-K) < 0$ (lihat \exerciseref{exercise:odddegnegativeK}).
Dalam pembuktiannya,
pastikan Anda menggunakan fakta bahwa $d$ ganjil.
Secara khusus, jika $d$ ganjil, maka ${(-n)}^d = -(n^d)$.

Kita menerapkan teorema nilai antara untuk memperoleh
$c \in (-K,M)$ sedemikian sehingga $g(c) = 0$.  Karena $g(x) = \frac{f(x)}{a_d}$,
maka $f(c) = 0$, dan pembuktian selesai.
\end{proof}

\begin{example}
Anda mungkin ingat betapa kerasnya kita bekerja
dalam \exampleref{example:sqrt2}
untuk menunjukkan bahwa $\sqrt{2}$ ada.
Dengan teorema Bolzano, kita dapat membuktikan keberadaan akar ke-$k$ dari setiap bilangan positif
$y > 0$ tanpa kesulitan.
Kita menyatakan bahwa untuk setiap $k \in \N$ dan setiap $y > 0$,
terdapat bilangan $x > 0$ sedemikian sehingga $x^k = y$.

Pembuktian: Jika $y=1$, pernyataannya jelas, jadi andaikan $y\neq 1$.
Demikian pula, andaikan $k \geq 2$.
Misalkan $f(x) \coloneqq x^k - y$.  Perhatikan bahwa $f(0) = -y < 0$.
Jika $y < 1$, maka $f(1) = 1^k -y > 0$.  Jika $y > 1$,
maka $f(y) = y^k-y = y(y^{k-1}-1) > 0$.
Dalam kedua kasus tersebut, terapkan teorema Bolzano untuk memperoleh $x > 0$ sedemikian sehingga $f(x) = 0$,
atau dengan kata lain $x^k = y$.
\end{example}

\begin{example}
Menariknya, terdapat fungsi-fungsi diskontinu yang memiliki
sifat nilai antara.
Fungsi
\begin{equation*}
f(x) \coloneqq
\begin{cases}
\sin(\nicefrac{1}{x}) & \text{jika } x \neq 0, \\
0 & \text{jika } x=0,
\end{cases}
\end{equation*}
tidak kontinu di $0$; namun, $f$ memiliki sifat nilai antara:
Setiap kali $a < b$ dan $y$ sedemikian sehingga $f(a) < y < f(b)$
atau $f(a) > y > f(b)$,
terdapat $c \in (a,b)$ sedemikian sehingga $f(c) = y$.
Lihat \figureref{figsin1x} untuk grafik $\sin(\nicefrac{1}{x})$.
Pembuktiannya diserahkan sebagai \exerciseref{exercise:meanvaluepropsin1x}.
\end{example}

Teorema nilai antara menyatakan bahwa jika $f \colon [a,b] \to \R$
kontinu, maka $f\bigl([a,b]\bigr)$ memuat semua nilai di antara $f(a)$ dan
$f(b)$.  Sebenarnya, pernyataan yang lebih kuat berlaku.  Dengan menggabungkan semua hasil dalam bagian ini,
kita dapat membuktikan akibat berguna berikut, yang pembuktiannya diserahkan sebagai latihan.
Petunjuk: Lihat \figureref{fig:imageinterval} dan perhatikan titik-titik ujung
interval citranya.

\begin{cor} \label{cor:imageofinterval}
Jika $f \colon [a,b] \to \R$ kontinu, maka citra langsung
$f\bigl([a,b]\bigr)$
merupakan interval tertutup dan terbatas atau suatu bilangan tunggal.
\end{cor}

\begin{myfigureht}
\myincludepdft{figimageinterval}{%
Grafik fungsi y sama dengan f dari x pada interval
dari a hingga b, yang ditandai dengan garis tebal.  Pada sumbu
vertikal, himpunan semua nilai yang dicapai oleh fungsi
ditandai sebagai f dari interval tertutup dari a hingga b.}
\caption{Citra fungsi kontinu $f \colon [a,b] \to \R$.\label{fig:imageinterval}}
\end{myfigureht}

\subsection{Latihan}

\begin{exercise}
Carilah contoh fungsi diskontinu $f \colon [0,1] \to \R$
yang menyebabkan kesimpulan teorema nilai antara gagal berlaku.
\end{exercise}

\begin{exercise}
Carilah contoh fungsi diskontinu \emph{terbatas} $f \colon [0,1]
\to \R$ yang tidak memiliki minimum mutlak maupun maksimum mutlak.
\end{exercise}

\begin{exercise}
Misalkan $f \colon (0,1) \to \R$ suatu fungsi kontinu sedemikian sehingga
$\displaystyle \lim_{x\to 0} f(x) =
\displaystyle \lim_{x\to 1} f(x) = 0$.  Buktikan bahwa
$f$ mencapai minimum mutlak atau maksimum mutlak pada $(0,1)$
(tetapi mungkin tidak keduanya).
\end{exercise}

\begin{exercise} \label{exercise:meanvaluepropsin1x}
Misalkan
\begin{equation*}
f(x) \coloneqq
\begin{cases}
\sin(\nicefrac{1}{x}) & \text{jika } x \neq 0, \\
0 & \text{jika } x=0.
\end{cases}
\end{equation*}
Buktikan bahwa $f$ memiliki sifat nilai antara.
Artinya, untuk setiap $a < b$, jika terdapat $y$ sedemikian sehingga $f(a) < y < f(b)$
atau $f(a) > y > f(b)$, maka
terdapat $c \in (a,b)$ sedemikian sehingga $f(c) = y$.
\end{exercise}

\begin{exercise} \label{exercise:odddegnegativeK}
Misalkan $g(x)$ suatu polinom monik berderajat ganjil $d$, yaitu,
\begin{equation*}
g(x) = x^d + b_{d-1} x^{d-1} + \cdots + b_1 x + b_0 ,
\end{equation*}
untuk beberapa bilangan real $b_{0}, b_1, \ldots, b_{d-1}$.  Buktikan bahwa terdapat
$K \in \N$ sedemikian sehingga $g(-K) < 0$.  Petunjuk: Pastikan Anda menggunakan fakta bahwa
$d$ ganjil.  Anda perlu menggunakan fakta bahwa ${(-n)}^d = -(n^d)$.
\end{exercise}

\begin{exercise}
Misalkan $g(x)$ suatu polinom monik berderajat genap positif $d$, yaitu,
\begin{equation*}
g(x) = x^d + b_{d-1} x^{d-1} + \cdots + b_1 x + b_0 ,
\end{equation*}
untuk beberapa bilangan real $b_{0}, b_1, \ldots, b_{d-1}$.  Andaikan
$g(0) < 0$.  Buktikan bahwa $g$ memiliki sekurang-kurangnya dua akar real yang berbeda.
\end{exercise}

\begin{exercise}
Buktikan \corref{cor:imageofinterval}:
Misalkan $f \colon [a,b] \to \R$ suatu fungsi kontinu.  Buktikan
bahwa citra langsung $f\bigl([a,b]\bigr)$ merupakan interval tertutup dan terbatas atau
sebuah bilangan tunggal.
\end{exercise}

\begin{exercise}
Misalkan $f \colon \R \to \R$ kontinu dan periodik dengan periode
$P > 0$.  Artinya, $f(x+P) = f(x)$ untuk setiap $x \in \R$.  Buktikan bahwa $f$
mencapai minimum mutlak dan maksimum mutlak.
\end{exercise}

\begin{exercise}[Menantang]
Misalkan $f(x)$ suatu polinom terbatas;
dengan kata lain, terdapat $M$ sedemikian sehingga $\babs{f(x)} \leq M$
untuk setiap $x \in \R$.  Buktikan bahwa $f$ harus konstan.
\end{exercise}

\begin{exercise}
Misalkan $f \colon [0,1] \to [0,1]$ kontinu.  Buktikan bahwa $f$
memiliki titik tetap; dengan kata lain, buktikan bahwa terdapat $x \in [0,1]$ sedemikian sehingga
$f(x) = x$.
\end{exercise}

\begin{exercise}
Carilah contoh fungsi kontinu terbatas $f \colon \R \to \R$ yang
tidak mencapai minimum mutlak maupun maksimum mutlak pada $\R$.
\end{exercise}

\begin{exercise}
Misalkan $f \colon \R \to \R$ kontinu sedemikian sehingga
$x \leq f(x) \leq x+1$ untuk setiap $x \in \R$.  Tentukan $f(\R)$.
\end{exercise}

\begin{exercise}
Benar atau salah? Buktikan atau carilah contoh tandingan.  Jika $f \colon \R \to
\R$ suatu fungsi kontinu sedemikian sehingga $f|_{\Z}$ terbatas, maka $f$
terbatas.
\end{exercise}

\begin{exercise}
Misalkan $f \colon [0,1] \to (0,1)$ suatu bijeksi.  Buktikan bahwa $f$ tidak
kontinu.
\end{exercise}

\begin{exercise}
Misalkan $f \colon \R \to \R$ kontinu.
\begin{enumerate}[a)]
\item
Buktikan bahwa jika terdapat $c$ sedemikian sehingga $f(c)f(-c) < 0$,
maka terdapat $d \in \R$ sedemikian sehingga $f(d) = 0$.
\item
Carilah fungsi kontinu $f$ sedemikian sehingga
$f(\R) = \R$, tetapi $f(x)f(-x) \geq 0$ untuk setiap $x \in \R$.
\end{enumerate}
\end{exercise}

\begin{exercise}
Misalkan $g(x)$ suatu polinom monik berderajat genap $d$, yaitu,
\begin{equation*}
g(x) = x^d + b_{d-1} x^{d-1} + \cdots + b_1 x + b_0 ,
\end{equation*}
untuk beberapa bilangan real $b_{0}, b_1, \ldots, b_{d-1}$.
Buktikan bahwa $g$ mencapai minimum mutlak pada $\R$.
\end{exercise}

\begin{exercise}
Misalkan $f(x)$ suatu polinom berderajat $d$ dan
$f(\R) = \R$.  Buktikan bahwa $d$ ganjil.
\end{exercise}

%%%%%%%%%%%%%%%%%%%%%%%%%%%%%%%%%%%%%%%%%%%%%%%%%%%%%%%%%%%%%%%%%%%%%%%%%%%%%%

\sectionnewpage
\section{Kekontinuan seragam}
\label{sec:unifcont}

%mbxINTROSUBSECTION

\sectionnotes{1,5--2 kuliah (perluasan kontinu dapat dijadikan materi opsional)}

\subsection{Kekontinuan seragam}

Kita sebelumnya menekankan bahwa $\delta$ dalam definisi kekontinuan
bergantung pada titik $c$.  Dalam beberapa situasi, kita diuntungkan jika dapat
memilih $\delta$ yang tidak bergantung pada titik mana pun; karena itu, konsep
ini kita beri nama.

\begin{defn}
Misalkan $S \subset \R$, dan misalkan $f \colon S \to \R$ suatu fungsi.
Andaikan untuk setiap $\epsilon > 0$ terdapat $\delta > 0$
sedemikian sehingga jika $x, c \in S$ dan
$\sabs{x-c} < \delta$, berlaku $\babs{f(x)-f(c)} < \epsilon$.
Kita katakan bahwa $f$ \emph{\myindex{kontinu seragam}}.
\end{defn}

Fungsi yang kontinu seragam tentu kontinu.
Satu-satunya perbedaan antara kedua definisi tersebut ialah bahwa dalam
kekontinuan seragam, untuk suatu $\epsilon > 0$ kita memilih $\delta > 0$
yang berlaku untuk semua $c \in S$.  Artinya, $\delta$ tidak lagi boleh
bergantung pada $c$; pilihan tersebut hanya bergantung pada $\epsilon$.  Kini domain
definisi fungsi berpengaruh.  Fungsi yang tidak kontinu seragam pada suatu
himpunan yang lebih besar dapat menjadi kontinu seragam ketika dibatasi pada
himpunan yang lebih kecil.
Perhatikan bahwa $x$ dan $c$ tidak diperlakukan secara berbeda dalam definisi
ini.

\begin{example}
$f \colon [0,1] \to \R$ yang didefinisikan oleh $f(x) \coloneqq x^2$ kontinu
seragam.

Bukti: Perhatikan bahwa $0 \leq x,c \leq 1$.  Maka
\begin{equation*}
\sabs{x^2-c^2} = \sabs{x+c}\sabs{x-c}
\leq \bigl(\sabs{x}+\sabs{c}\bigr) \sabs{x-c}
\leq (1+1)\sabs{x-c} .
\end{equation*}
Jadi, untuk $\epsilon > 0$, ambil $\delta \coloneqq \nicefrac{\epsilon}{2}$.
Jika $\sabs{x-c} < \delta$, maka $\sabs{x^2-c^2} \leq 2 \sabs{x-c} < \epsilon$.

\medskip

Di sisi lain, $g \colon \R \to \R$ yang didefinisikan oleh
$g(x) \coloneqq x^2$ tidak kontinu seragam.

Bukti: Andaikan fungsi tersebut kontinu seragam.
Maka, untuk setiap $\epsilon > 0$,
akan terdapat $\delta > 0$ sedemikian sehingga
jika $\sabs{x-c} < \delta$, maka $\sabs{x^2 -c^2} < \epsilon$.
Ambil $x > 0$ dan tetapkan
$c \coloneqq x+\nicefrac{\delta}{2}$.  Kita peroleh
\begin{equation*}
\epsilon >
\sabs{x^2-c^2} = \sabs{x+c}\sabs{x-c}
=
(2x+\nicefrac{\delta}{2})\nicefrac{\delta}{2} 
\geq 
\delta x .
\end{equation*}
Jadi, $x < \nicefrac{\epsilon}{\delta}$ untuk semua $x > 0$, suatu
kontradiksi.
\end{example}


\begin{example}
Fungsi $f \colon (0,1) \to \R$ yang didefinisikan oleh
$f(x) \coloneqq \nicefrac{1}{x}$ tidak kontinu seragam.

Bukti: Untuk $\epsilon > 0$, ketaksamaan $\epsilon >
\abs{\nicefrac{1}{x}-\nicefrac{1}{y}}$ berlaku jika dan hanya jika
\begin{equation*}
\epsilon >
\abs{\nicefrac{1}{x}-\nicefrac{1}{y}}
=
\frac{\sabs{y-x}}{\sabs{xy}} 
=
\frac{\sabs{y-x}}{xy} ,
\end{equation*}
atau
\begin{equation*}
\sabs{x-y} < xy \epsilon .
\end{equation*}
Andaikan $\epsilon < 1$.
Kita ingin mengetahui apakah suatu $\delta > 0$ yang kecil akan berlaku.
Jika $x \in (0,1)$ dan
$y = x+\nicefrac{\delta}{2} \in (0,1)$,
maka $\sabs{x-y} = \nicefrac{\delta}{2} < \delta$.
Dengan menyubstitusikan $y$ ke dalam ketaksamaan di atas, kita memperoleh
$\nicefrac{\delta}{2} <
x \bigl( x+\nicefrac{\delta}{2} \bigr) \epsilon < x$.
Jika definisi kekontinuan seragam terpenuhi, maka ketaksamaan
$\nicefrac{\delta}{2} < x$ akan berlaku untuk semua $x > 0$ yang kecil.
Namun, hal itu mengakibatkan $\delta \leq 0$.
Jadi, tidak ada satu pun $\delta > 0$ yang berlaku untuk semua titik.
\end{example}

Contoh-contoh tersebut menunjukkan bahwa jika $f$ didefinisikan pada interval
yang tidak tertutup atau tidak terbatas, maka $f$ mungkin kontinu tetapi tidak
kontinu seragam.  Namun, untuk interval tertutup dan terbatas $[a,b]$, kita
dapat menyatakan hasil berikut.

\begin{thm} \label{unifcont:thm}
Misalkan $f \colon [a,b] \to \R$ suatu fungsi kontinu.  Maka $f$
kontinu seragam.
\end{thm}

\begin{proof}
Kita membuktikan pernyataan tersebut dengan kontraposisi.
Andaikan $f$ tidak kontinu seragam.  Kita akan membuktikan bahwa terdapat
$c \in [a,b]$ sedemikian sehingga $f$ tidak kontinu di titik itu.  Negasi definisi kekontinuan seragam
adalah sebagai berikut: Terdapat $\epsilon > 0$ sedemikian sehingga untuk
setiap $\delta > 0$, terdapat titik-titik $x, y$ dalam $[a,b]$ dengan
$\sabs{x-y} < \delta$ dan $\babs{f(x)-f(y)} \geq \epsilon$.

Jadi, untuk $\epsilon > 0$ di atas, kita memperoleh barisan
$\{ x_n \}_{n=1}^\infty$ dan $\{ y_n \}_{n=1}^\infty$ sedemikian sehingga
$\sabs{x_n-y_n} < \nicefrac{1}{n}$ serta $\babs{f(x_n)-f(y_n)} \geq
\epsilon$.  Berdasarkan
\hyperref[thm:bwseq]{Bolzano--Weierstrass},
terdapat subbarisan konvergen
$\{ x_{n_k} \}_{k=1}^\infty$.  Tetapkan $c \coloneqq \lim_{k\to\infty} x_{n_k}$.
Karena $a \leq x_{n_k} \leq b$ untuk semua $k$, kita memperoleh
$a \leq c \leq b$.  Kita estimasikan
\begin{equation*}
\sabs{y_{n_k} - c} =
\sabs{y_{n_k} - x_{n_k} + x_{n_k} - c} \leq
\sabs{y_{n_k} - x_{n_k}}
+
\sabs{x_{n_k}-c}
<
\nicefrac{1}{n_k} 
+
\sabs{x_{n_k}-c} .
\end{equation*}
Karena $\nicefrac{1}{n_k}$ dan $\sabs{x_{n_k}-c}$ keduanya menuju nol ketika
$k$ menuju tak hingga, $\{ y_{n_k} \}_{k=1}^\infty$ konvergen dan limitnya
adalah $c$.  Sekarang kita tunjukkan bahwa $f$ tidak kontinu di $c$.
Kita estimasikan
\begin{equation*}
\begin{split}
\babs{f(x_{n_k}) - f(c)} & =
\babs{f(x_{n_k}) - f(y_{n_k}) + f(y_{n_k}) - f(c)} \\
& \geq
\babs{f(x_{n_k}) - f(y_{n_k})} - \babs{f(y_{n_k}) - f(c)} \\
& \geq
\epsilon - \babs{f(y_{n_k})-f(c)} .
\end{split}
\end{equation*}
Dengan kata lain,
\begin{equation*}
\babs{f(x_{n_k})-f(c)} 
+
\babs{f(y_{n_k})-f(c)}  \geq
\epsilon .
\end{equation*}
Setidaknya salah satu barisan $\bigl\{ f(x_{n_k}) \bigr\}_{k=1}^\infty$ atau
$\bigl\{ f(y_{n_k}) \bigr\}_{k=1}^\infty$ tidak dapat konvergen ke $f(c)$.
Jika tidak, ruas kiri ketaksamaan tersebut akan menuju nol, sedangkan ruas
kanannya positif.  Dengan demikian, $f$ tidak mungkin kontinu di $c$.
\end{proof}

Seperti sebelumnya, perhatikan unsur kunci dalam pembuktian: Kita dapat
menerapkan Bolzano--Weierstrass karena interval $[a,b]$ terbatas, dan limit
subbarisan kembali berada dalam $[a,b]$ karena interval tersebut tertutup.

\subsection{Perluasan kontinu}

Fungsi yang kontinu seragam pada interval terbuka dapat diperluas secara kontinu
ke titik-titik ujungnya.
Kuncinya adalah lemma berikut, yang juga memiliki banyak kegunaan lain.
Lemma tersebut mengatakan bahwa fungsi kontinu seragam mempertahankan barisan Cauchy.
Persoalan baru di sini adalah bahwa untuk suatu barisan Cauchy,
limitnya mungkin tidak berada dalam domain fungsi.

\begin{lemma} \label{unifcauchycauchy:lemma}
Misalkan $S \subset \R$ dan misalkan $f \colon S \to \R$ suatu fungsi kontinu seragam.  Misalkan
$\{ x_n \}_{n=1}^\infty$ suatu barisan Cauchy dalam $S$.
Maka $\bigl\{ f(x_n) \bigr\}_{n=1}^\infty$ merupakan barisan Cauchy.
\end{lemma}

\begin{proof}
Ambil sebarang $\epsilon > 0$.  Terdapat $\delta > 0$ sedemikian sehingga
$\babs{f(x)-f(y)} < \epsilon$ setiap kali $x,y \in S$ dan $\sabs{x-y} < \delta$.  Pilih $M
\in \N$ sedemikian sehingga untuk semua $n, k \geq M$, berlaku $\sabs{x_n-x_k} < \delta$.
Maka untuk semua $n, k \geq M$, berlaku $\babs{f(x_n)-f(x_k)} < \epsilon$.
\end{proof}

Salah satu penerapan lemma di atas adalah hasil perluasan berikut.  Hasil ini mengatakan bahwa
suatu fungsi pada interval terbuka kontinu seragam jika dan hanya jika
fungsi tersebut dapat diperluas menjadi fungsi kontinu pada interval tertutup.

\begin{prop} \label{context:prop}
Fungsi $f \colon (a,b) \to \R$ kontinu seragam jika dan hanya jika
kedua limit
\begin{equation*}
L_a \coloneqq \lim_{x \to a} f(x) \qquad \text{dan} \qquad
L_b \coloneqq \lim_{x \to b} f(x)
\end{equation*}
ada dan fungsi $\widetilde{f} \colon [a,b] \to \R$
yang didefinisikan oleh
\begin{equation*}
\widetilde{f}(x) \coloneqq
\begin{cases}
f(x) & \text{jika } x \in (a,b), \\
L_a & \text{jika } x = a, \\
L_b & \text{jika } x = b
\end{cases}
\end{equation*}
bersifat kontinu.
\end{prop}

\begin{proof}
Salah satu arahnya singkat.  Jika $\widetilde{f}$ kontinu, maka
fungsi tersebut kontinu seragam berdasarkan \thmref{unifcont:thm}.  Karena $f$ merupakan
pembatasan $\widetilde{f}$ pada $(a,b)$, fungsi $f$ juga kontinu seragam
(latihan).

Sekarang andaikan $f$ kontinu seragam.  Pertama-tama kita harus menunjukkan
bahwa limit $L_a$ dan $L_b$ ada.  Mari kita pusatkan perhatian pada $L_a$.
Ambil $\{ x_n \}_{n=1}^\infty$ dalam $(a,b)$ sedemikian sehingga
$\lim_{n\to\infty} x_n = a$.
Barisan $\{ x_n \}_{n=1}^\infty$ merupakan barisan Cauchy, sehingga berdasarkan
\lemmaref{unifcauchycauchy:lemma}
barisan $\bigl\{ f(x_n) \bigr\}_{n=1}^\infty$ merupakan barisan Cauchy dan karenanya konvergen.
Misalkan $L_1 \coloneqq \lim_{n\to\infty} f(x_n)$.  Ambil barisan lain
$\{ y_n \}_{n=1}^\infty$ dalam $(a,b)$ sedemikian sehingga $\lim_{n\to\infty} y_n = a$.
Dengan penalaran yang sama,
kita peroleh $L_2 \coloneqq \lim_{n\to\infty} f(y_n)$.  Jika kita menunjukkan bahwa $L_1 = L_2$, maka
limit $L_a = \lim_{x\to a} f(x)$ ada.  Ambil sebarang $\epsilon > 0$.
Pilih $\delta > 0$ sedemikian sehingga $\sabs{x-y} < \delta$ mengakibatkan $\babs{f(x)-f(y)} <
\nicefrac{\epsilon}{3}$.  Pilih $M \in \N$ sedemikian sehingga untuk
$n \geq M$, berlaku $\sabs{a-x_n} < \nicefrac{\delta}{2}$,
$\sabs{a-y_n} < \nicefrac{\delta}{2}$,
$\babs{f(x_n)-L_1} < \nicefrac{\epsilon}{3}$, dan
$\babs{f(y_n)-L_2} < \nicefrac{\epsilon}{3}$.  Maka untuk $n \geq M$,
\begin{equation*}
\sabs{x_n-y_n} = 
\sabs{x_n-a+a-y_n} \leq
\sabs{x_n-a}+\sabs{a-y_n} < \nicefrac{\delta}{2} + \nicefrac{\delta}{2} =
\delta.
\end{equation*}
Jadi,
\begin{equation*}
\begin{split}
\sabs{L_1-L_2} &=
\babs{L_1-f(x_n)+f(x_n)-f(y_n)+f(y_n)-L_2} \\
& \leq 
\babs{L_1-f(x_n)}+\babs{f(x_n)-f(y_n)}+\babs{f(y_n)-L_2} \\
& <
\nicefrac{\epsilon}{3} + \nicefrac{\epsilon}{3} + \nicefrac{\epsilon}{3}
=
\epsilon .
\end{split}
\end{equation*}
Oleh karena itu, $L_1 = L_2$.
Jadi, $L_a$ ada.  Pembuktian bahwa $L_b$ ada diserahkan sebagai latihan.

Jika $L_a = \lim_{x\to a} f(x)$
ada, maka $\lim_{x\to a} \widetilde{f}(x)$ ada
dan sama dengan $L_a$
(lihat \propref{prop:limrest}).  Hal yang sama berlaku untuk $L_b$.
Jadi, $\widetilde{f}$ kontinu di $a$ dan $b$.
Karena $f$ kontinu di $c \in (a,b)$,
kita peroleh bahwa
$\widetilde{f}$ kontinu di $c \in (a,b)$ (gunakan kembali \propref{prop:limrest}).
\end{proof}

Penerapan umum proposisi ini (bersama
\propref{prop:onesidedlimits})
adalah sebagai berikut.
Jika $f \colon (-1,0) \cup (0,1) \to \R$ kontinu seragam,
maka $\lim_{x\to 0} f(x)$ ada dan fungsi tersebut
memiliki \emph{\myindex{singularitas yang dapat dihilangkan}}.
Artinya,
kita dapat memperluas fungsi tersebut menjadi fungsi kontinu pada $(-1,1)$.


\subsection{Fungsi kontinu Lipschitz}

\begin{defn}
Suatu fungsi $f \colon S \to \R$
disebut \emph{\myindex{kontinu Lipschitz}}\index{fungsi!Lipschitz}%
\footnote{Dinamai menurut matematikawan Jerman
\href{https://en.wikipedia.org/wiki/Rudolf_Lipschitz}{Rudolf Otto Sigismund Lipschitz}
(1832--1903).} jika terdapat $K \in \R$ sedemikian sehingga
\begin{equation*}
\babs{f(x)-f(y)} \leq K \sabs{x-y} 
\qquad \text{untuk semua } x \text{ dan } y \text{ dalam } S.
\end{equation*}
\end{defn}

Kelas fungsi yang bersifat kontinu Lipschitz sangat luas.  Namun, berhati-hatilah.
Seperti halnya untuk fungsi kontinu seragam,
domain definisi fungsi tersebut penting.  Lihat contoh-contoh di bawah
dan bagian latihan.  Pertama, kita benarkan penggunaan kata \emph{kontinu}.

\begin{prop}
Fungsi kontinu Lipschitz bersifat kontinu seragam.
\end{prop}

\begin{proof}
Misalkan $f \colon S \to \R$ suatu fungsi dan pilih konstanta positif $K$ sedemikian sehingga
$\babs{f(x)-f(y)} \leq K \sabs{x-y}$
untuk semua $x, y$ dalam $S$.
Ambil sebarang $\epsilon > 0$.  Pilih $\delta \coloneqq
\nicefrac{\epsilon}{K}$.
Untuk semua $x$ dan $y$ dalam $S$ sedemikian sehingga
$\sabs{x-y} < \delta$,
\begin{equation*}
\babs{f(x)-f(y)} \leq K \sabs{x-y} < K \delta = K \frac{\epsilon}{K} =
\epsilon .
\end{equation*}
Oleh karena itu, $f$ kontinu seragam.
\end{proof}

Kita tafsirkan kekontinuan Lipschitz secara geometris.  Misalkan $f$ suatu fungsi
kontinu Lipschitz dengan suatu konstanta $K$.  Ketaksamaan tersebut dapat ditulis ulang
sehingga untuk $x \neq y$ berlaku
\begin{equation*}
\abs{\frac{f(x)-f(y)}{x-y}} \leq K .
\end{equation*}
Besaran $\frac{f(x)-f(y)}{x-y}$ merupakan kemiringan garis
yang melalui titik $\bigl(x,f(x)\bigr)$
dan $\bigl(y,f(y)\bigr)$, yaitu suatu \emph{\myindex{garis sekan}}.  Oleh karena itu, $f$
kontinu Lipschitz jika dan hanya jika setiap garis yang memotong grafik $f$ di setidaknya dua
titik berbeda memiliki nilai mutlak kemiringan yang kurang dari atau sama dengan $K$.
Lihat \figureref{fig:lipschitz}.
\begin{myfigureht}
\myincludegraphics{lipschitzfig}{%
Grafik suatu fungsi dengan beberapa titik sudut.  Dua
titik pada sumbu horizontal ditandai dengan x
dan y.  Sebuah garis digambar melalui titik-titik yang bersesuaian
pada grafik, dan kemiringannya diberi label
f dari x dikurangi f dari y, semuanya dibagi x dikurangi y.}
\caption{Kemiringan suatu garis sekan.
Suatu fungsi bersifat Lipschitz jika %$\abs{\text{slope}} =
$\abs{\frac{f(x)-f(y)}{x-y}} \leq K$ untuk semua $x$ dan $y$.\label{fig:lipschitz}}
\end{myfigureht}

\begin{example}
Fungsi $\sin(x)$ dan $\cos(x)$ bersifat kontinu Lipschitz.
Dalam \exampleref{sincos:example} kita telah melihat kedua ketaksamaan berikut.
\begin{equation*}
\babs{\sin(x)-\sin(y)} 
\leq \sabs{x-y}
\qquad \text{dan} \qquad
\babs{\cos(x)-\cos(y)}
\leq \sabs{x-y} .
\end{equation*}

Jadi, sinus dan kosinus bersifat kontinu Lipschitz dengan $K=1$.
\end{example}

\begin{example}
Fungsi $f \colon [1,\infty) \to \R$ yang didefinisikan oleh $f(x) \coloneqq \sqrt{x}$
bersifat kontinu Lipschitz. Buktinya:
\begin{equation*}
\abs{\sqrt{x}-\sqrt{y}} = 
\abs{\frac{x-y}{\sqrt{x}+\sqrt{y}}}
=
\frac{\sabs{x-y}}{\sqrt{x}+\sqrt{y}} .
\end{equation*}
Karena $x \geq 1$ dan $y \geq 1$, kita peroleh $\frac{1}{\sqrt{x}+\sqrt{y}}
\leq \frac{1}{2}$.  Oleh karena itu,
\begin{equation*}
\abs{\sqrt{x}-\sqrt{y}} = 
\abs{\frac{x-y}{\sqrt{x}+\sqrt{y}}}
\leq
\frac{1}{2}
\sabs{x-y}.
\end{equation*}

Di sisi lain, $g \colon [0,\infty) \to \R$ yang didefinisikan oleh
$g(x) \coloneqq \sqrt{x}$ tidak kontinu Lipschitz.  Buktinya:
Andaikan terdapat $K$ sedemikian sehingga untuk semua $x,y \in [0,\infty)$,
\begin{equation*}
\abs{\sqrt{x}-\sqrt{y}} 
\leq
K \sabs{x-y} ,
\end{equation*}
Kita peroleh $K > 0$, sebab $K < 0$ mustahil,
dan $K=0$ akan berarti bahwa $g$ konstan.
Tetapkan $y=0$ untuk memperoleh
$\sqrt{x} \leq K x$.
Untuk $x > 0$, kita peroleh
$\nicefrac{1}{K} \leq \sqrt{x}$ atau $\nicefrac{1}{K^2} \leq x$.
Hal ini tidak mungkin benar untuk semua
$x > 0$.  Jadi, tidak ada $K$ yang memenuhi dan $g$ tidak
kontinu Lipschitz.  Lihat \figureref{fig:sqrtgraph} dan perhatikan bahwa garis-garis
sekan semakin tegak ketika kita mendekati $x=0$.
\begin{myfigureht}
\myincludegraphics{sqrtgraph}{%
Grafik suatu fungsi yang mula-mula naik secara vertikal
di titik ujung kiri, lalu cepat melengkung sehingga menjadi
jauh lebih mendatar di titik ujung kanan.
Digambar garis-garis sekan yang melalui titik ujung kiri dan suatu
barisan titik pada grafik yang semakin mendekati
titik ujung kiri.
Tampak bahwa garis-garis sekan semakin
tegak ketika titik kedua
pada garis sekan semakin dekat
ke titik ujung kiri.}
\caption{Grafik $\sqrt{x}$ dan beberapa garis sekan yang melalui $(0,0)$ dan
$(x,\sqrt{x})$.\label{fig:sqrtgraph}}
\end{myfigureht}

Fungsi $g$ pada contoh terakhir kontinu seragam,
tetapi tidak kontinu Lipschitz.
Untuk melihat bahwa $\sqrt{x}$ merupakan fungsi yang
kontinu seragam pada $[0,\infty)$,
perhatikan bahwa pembatasannya pada $[0,1]$ kontinu seragam berdasarkan \thmref{unifcont:thm}.
Pembatasannya pada $[1,\infty)$ juga bersifat Lipschitz (dan karenanya kontinu seragam).
Tidak sulit (latihan) untuk menunjukkan bahwa ini berarti $\sqrt{x}$ merupakan
fungsi kontinu seragam pada $[0,\infty)$.
\end{example}

\subsection{Latihan}

\begin{exercise}
Misalkan $f \colon S \to \R$ kontinu seragam.  Misalkan $A \subset S$.
Maka pembatasan $f|_A$ kontinu seragam.
\end{exercise}

\begin{exercise}
Misalkan $f \colon (a,b) \to \R$ suatu fungsi kontinu seragam.
Selesaikan pembuktian \propref{context:prop} dengan menunjukkan bahwa
limit
$\lim\limits_{x \to b} f(x)$
ada.
\end{exercise}

\begin{exercise}
Tunjukkan bahwa $f \colon (c,\infty) \to \R$ untuk suatu $c > 0$
yang didefinisikan oleh $f(x) \coloneqq \nicefrac{1}{x}$ bersifat kontinu Lipschitz.
\end{exercise}

\begin{exercise}
Tunjukkan bahwa $f \colon (0,\infty) \to \R$
yang didefinisikan oleh $f(x) \coloneqq \nicefrac{1}{x}$ tidak kontinu Lipschitz.
\end{exercise}

\begin{samepage}
\begin{exercise}
Misalkan $A$ dan $B$ merupakan interval.
Misalkan $f \colon A \to \R$ dan $g \colon B \to \R$ fungsi-fungsi kontinu seragam
sedemikian sehingga $f(x) = g(x)$ untuk $x \in A \cap B$.  Definisikan
fungsi $h \colon A \cup B \to \R$ dengan $h(x) \coloneqq f(x)$ jika
$x \in A$ dan $h(x) \coloneqq g(x)$ jika $x \in B \setminus A$.
\begin{enumerate}[a)]
\item
Buktikan bahwa jika $A \cap B \neq \emptyset$, maka $h$ kontinu seragam.
\item
Temukan contoh dengan $A \cap B = \emptyset$ dan $h$ bahkan tidak
kontinu.
\end{enumerate}
\end{exercise}
\end{samepage}

\begin{exercise}[Menantang]
Misalkan $f \colon \R \to \R$ suatu polinomial berderajat
$d \geq 2$.  Tunjukkan bahwa $f$ tidak kontinu
Lipschitz.
\end{exercise}

\begin{exercise}
Misalkan $f \colon (0,1) \to \R$ suatu fungsi kontinu terbatas.  Tunjukkan bahwa
fungsi
$g(x) \coloneqq x(1-x)f(x)$ kontinu seragam.
\end{exercise}

\begin{exercise}
Tunjukkan bahwa $f \colon (0,\infty) \to \R$ yang didefinisikan oleh $f(x) \coloneqq \sin
(\nicefrac{1}{x})$ tidak kontinu seragam.
\end{exercise}

\begin{exercise}[Menantang]
Misalkan $f \colon \Q \to \R$ suatu fungsi kontinu seragam.  Tunjukkan bahwa
terdapat fungsi kontinu seragam $\widetilde{f} \colon \R \to \R$
sedemikian sehingga $f(x) = \widetilde{f}(x)$ untuk semua $x \in \Q$.
\end{exercise}

\begin{samepage}
\begin{exercise}
\leavevmode
\begin{enumerate}[a)]
\item
Temukan fungsi kontinu $f \colon (0,1) \to \R$ dan barisan
$\{ x_n \}_{n=1}^\infty$ dalam
$(0,1)$ yang merupakan barisan Cauchy, tetapi $\bigl\{ f(x_n) \bigr\}_{n=1}^\infty$ bukan barisan Cauchy.
\item
Buktikan bahwa jika $f \colon \R \to \R$ kontinu,
dan $\{ x_n \}_{n=1}^\infty$ merupakan
barisan Cauchy, maka $\bigl\{ f(x_n) \bigr\}_{n=1}^\infty$ juga merupakan barisan Cauchy.
\end{enumerate}
\end{exercise}
\end{samepage}

\begin{samepage}
\begin{exercise}
Buktikan:
\begin{enumerate}[a)]
\item
Jika $f \colon S \to \R$ dan $g \colon S \to \R$ kontinu seragam,
maka $h \colon S \to \R$ yang diberikan oleh $h(x) \coloneqq f(x) + g(x)$
kontinu seragam.
\item
Jika $f \colon S \to \R$ kontinu seragam dan $a \in \R$,
maka $h \colon S \to \R$ yang diberikan oleh $h(x) \coloneqq a f(x)$
kontinu seragam.
\end{enumerate}
\end{exercise}
\end{samepage}

\begin{exercise}
Buktikan:
\begin{enumerate}[a)]
\item
Jika $f \colon S \to \R$ dan $g \colon S \to \R$ bersifat Lipschitz,
maka $h \colon S \to \R$ yang diberikan oleh $h(x) \coloneqq f(x) + g(x)$
bersifat Lipschitz.
\item
Jika $f \colon S \to \R$ bersifat Lipschitz dan $a \in \R$,
maka $h \colon S \to \R$ yang diberikan oleh $h(x) \coloneqq a f(x)$
bersifat Lipschitz.
\end{enumerate}
\end{exercise}

\begin{exercise}
\pagebreak[2]
\leavevmode
\begin{enumerate}[a)]
\item
Jika $f \colon [0,1] \to \R$ diberikan oleh $f(x) \coloneqq x^m$ untuk suatu bilangan bulat
$m \geq 0$,
tunjukkan bahwa $f$ bersifat Lipschitz dan temukan konstanta Lipschitz terbaik (terkecil)
$K$ (yang tentu saja bergantung pada $m$).
Petunjuk:
$(x-y)(x^{m-1} + x^{m-2}y + x^{m-3}y^2 + \cdots + x y^{m-2} + y^{m-1}) = x^m
- y^m$ untuk $m$ positif.
\item
Dengan menggunakan latihan sebelumnya, tunjukkan bahwa jika $f \colon [0,1] \to \R$
suatu polinomial, yaitu $f(x) \coloneqq a_m x^m + a_{m-1} x^{m-1} + \cdots + a_0$,
maka $f$ bersifat Lipschitz.
\end{enumerate}
\end{exercise}

\begin{exercise}
\pagebreak[2]
Andaikan untuk $f \colon [0,1] \to \R$ berlaku $\babs{f(x)-f(y)} \leq K
\sabs{x-y}$ untuk semua $x,y$ dalam $[0,1]$,
dan $f(0) = f(1) = 0$.
Buktikan bahwa $\babs{f(x)} \leq \nicefrac{K}{2}$ untuk semua $x \in [0,1]$.  Tunjukkan pula dengan contoh bahwa
$\nicefrac{K}{2}$ merupakan batas terbaik yang mungkin.
Artinya, temukan fungsi kontinu
yang memenuhi $\babs{f(x)} = \nicefrac{K}{2}$ untuk suatu $x \in [0,1]$.
\end{exercise}

\begin{exercise}
Andaikan $f \colon \R \to \R$ kontinu dan periodik dengan periode
$P > 0$.  Artinya, $f(x+P) = f(x)$ untuk semua $x \in \R$.  Tunjukkan bahwa $f$
kontinu seragam.
\end{exercise}

\begin{exercise}
Andaikan $f \colon S \to \R$ dan $g \colon [0,\infty) \to [0,\infty)$
merupakan fungsi, $g$ kontinu di $0$, $g(0) = 0$, dan
setiap kali $x$ dan $y$ berada dalam $S$, berlaku $\babs{f(x)-f(y)} \leq
g\bigl(\sabs{x-y}\bigr)$.
Buktikan bahwa $f$ kontinu seragam.
\end{exercise}

\begin{exercise}
Andaikan $f \colon [a,b] \to \R$ suatu fungsi sedemikian sehingga untuk setiap $c \in
[a,b]$ terdapat $K_c > 0$ dan $\epsilon_c > 0$ yang memenuhi
$\babs{f(x)-f(y)} \leq K_c \sabs{x-y}$ untuk semua $x$ dan $y$ dalam
$(c-\epsilon_c,c+\epsilon_c) \cap [a,b]$.  Dengan kata lain, $f$
\myquote{bersifat Lipschitz lokal.}
\begin{enumerate}[a)]
\item
Buktikan bahwa terdapat satu $K > 0$ sedemikian sehingga
$\babs{f(x)-f(y)} \leq K \sabs{x-y}$ untuk semua $x,y$ dalam $[a,b]$.
\item
Temukan contoh tandingan terhadap pernyataan di atas jika intervalnya terbuka, yaitu,
temukan $f \colon (a,b) \to \R$ yang bersifat Lipschitz lokal, tetapi tidak
Lipschitz.
\end{enumerate}
\end{exercise}

%%%%%%%%%%%%%%%%%%%%%%%%%%%%%%%%%%%%%%%%%%%%%%%%%%%%%%%%%%%%%%%%%%%%%%%%%%%%%%

\sectionnewpage
\section{Limit di tak hingga}
\label{sec:limitatinf}

%mbxINTROSUBSECTION

\sectionnotes{kurang dari 1 kuliah (opsional, dapat dihilangkan dengan aman kecuali
\sectionref{sec:monotonefunc} atau
\sectionref{sec:impropriemann} juga dibahas)}

\subsection{Limit di tak hingga}

Seperti halnya barisan, suatu variabel kontinu juga dapat menuju tak hingga.

\begin{defn}
Kita mengatakan bahwa $\infty$ merupakan \emph{\myindex{titik akumulasi}} dari $S \subset \R$ jika untuk setiap
$M \in \R$, terdapat $x \in S$ sedemikian sehingga $x \geq M$.  Demikian pula,
$- \infty$ merupakan \emph{titik akumulasi} dari $S \subset \R$ jika untuk setiap
$M \in \R$, terdapat $x \in S$ sedemikian sehingga $x \leq M$.

\index{limit!fungsi di tak hingga}%
Misalkan $f \colon S \to \R$ suatu fungsi, dengan 
$\infty$ merupakan titik akumulasi dari $S$.
Jika terdapat $L \in \R$
dengan sifat
bahwa untuk setiap $\epsilon > 0$, terdapat $M \in \R$ sedemikian sehingga
\begin{equation*}
\babs{f(x) - L} < \epsilon 
\end{equation*}
setiap kali $x \in S$ dan $x \geq M$, kita mengatakan bahwa $f(x)$
\emph{konvergen}\index{konvergen!fungsi} ke $L$
ketika $x$ menuju $\infty$.  Kita menyebut $L$ suatu \emph{limit}
dan, jika tunggal, menuliskan
\glsadd{not:limfuncinf}
\begin{equation*}
\lim_{x \to \infty} f(x) \coloneqq L .
\end{equation*}
Sebagai alternatif, kita menuliskan $f(x) \to L$ ketika $x \to \infty$.

Demikian pula, jika $-\infty$ merupakan titik akumulasi dari $S$
dan
terdapat $L \in \R$
sedemikian sehingga untuk setiap $\epsilon > 0$, terdapat $M \in \R$ sedemikian sehingga
\begin{equation*}
\babs{f(x) - L} < \epsilon 
\end{equation*}
setiap kali $x \in S$ dan $x \leq M$, kita mengatakan bahwa $f(x)$ \emph{konvergen} ke $L$
ketika $x$ menuju $-\infty$.
Sebagai alternatif, kita menuliskan $f(x) \to L$ ketika $x \to -\infty$.
Kita menyebut $L$ suatu \emph{limit} dan, jika tunggal, menuliskan
\begin{equation*}
\lim_{x \to -\infty} f(x) \coloneqq L .
\end{equation*}
\end{defn}

Hal pertama yang perlu dilakukan, seperti biasa, adalah membuktikan bahwa limit tersebut, jika ada,
bersifat tunggal.
Kita menyerahkannya sebagai latihan bagi pembaca.

\begin{prop} \label{liminfty:unique}
Limit di $\infty$ atau $-\infty$ sebagaimana didefinisikan di atas bersifat tunggal jika ada.
\end{prop}

\begin{example}
Misalkan $f(x) \coloneqq \frac{1}{\sabs{x}+1}$.  Maka
\begin{equation*}
\lim_{x\to \infty} f(x) = 0 \qquad \text{dan} \qquad
\lim_{x\to -\infty} f(x) = 0 .
\end{equation*}

Bukti:
Ambil $\epsilon > 0$.  Cari $M > 0$ yang cukup besar
sehingga $\frac{1}{M+1} < \epsilon$.  Jika
$x \geq M$, maka $0 < \frac{1}{\sabs{x}+1} = \frac{1}{x+1} \leq \frac{1}{M+1} < \epsilon$.
Limit pertama pun diperoleh.
Bukti untuk $-\infty$ diserahkan kepada pembaca.
\end{example}

\begin{example}
Misalkan $f(x) \coloneqq \sin(\pi x)$.  Maka $\lim_{x\to\infty} f(x)$ tidak ada.
Untuk membuktikan fakta ini, perhatikan bahwa jika $x = 2n+\nicefrac{1}{2}$ untuk suatu $n \in
\N$, maka $f(x)=1$,
sedangkan jika $x = 2n+\nicefrac{3}{2}$, maka $f(x)=-1$.
Kedua nilai ini, $1$ dan $-1$, tidak mungkin sekaligus berada
dalam jarak $\epsilon$ yang kecil dari satu bilangan real yang sama.

Berhati-hatilah agar tidak mencampuradukkan limit kontinu dengan limit barisan.
Kita dapat mengatakan
\begin{equation*}
\lim_{n \to \infty} \sin(\pi n) = 0, \qquad \text{tetapi} \qquad
\lim_{x \to \infty} \sin(\pi x) \enspace \text{tidak ada}.
\end{equation*}
Tentu saja, notasi tersebut ambigu:  Apakah yang kita maksud adalah
barisan $\bigl\{ \sin (\pi n) \bigr\}_{n=1}^\infty$ atau fungsi $\sin(\pi x)$
dari suatu variabel real?  Kita sekadar menggunakan konvensi
bahwa $n \in \N$, sedangkan $x \in \R$.  Ketika notasinya tidak jelas,
sebaiknya nyatakan secara eksplisit di himpunan mana variabel tersebut berada, atau jenis
limit apa yang sedang digunakan.  Jika ada kemungkinan timbul kebingungan, kita dapat
menuliskan, misalnya,
\begin{equation*}
\lim_{\substack{n \to \infty\\n \in \N}} \sin(\pi n) .
\end{equation*}
\end{example}

Terdapat hubungan antara limit kontinu dan limit barisan, tetapi kita harus meninjau semua
barisan yang menuju tak hingga, seperti sebelumnya dalam \lemmaref{seqflimit:lemma}.

\begin{lemma} \label{seqflimitinf:lemma}
Misalkan $f \colon S \to \R$ suatu fungsi, $\infty$ suatu titik
akumulasi dari $S \subset \R$, dan $L \in \R$.  Maka
\begin{equation*}
\lim_{x\to\infty} f(x) = L
\qquad \text{jika dan hanya jika} \qquad
\lim_{n\to\infty} f(x_n) = L
\end{equation*}
untuk setiap barisan $\{ x_n \}_{n=1}^\infty$ dalam $S$ sedemikian sehingga $\lim\limits_{n\to\infty} x_n = \infty$.
\end{lemma}

Lemma tersebut juga berlaku untuk limit ketika $x \to -\infty$.
Buktinya hampir identik dan
diserahkan sebagai latihan.

\begin{proof}
Pertama, misalkan $f(x) \to L$ ketika $x \to \infty$.
Untuk sebarang $\epsilon > 0$, terdapat $M$ sedemikian sehingga untuk semua $x \in S$
yang memenuhi $x \geq M$,
berlaku $\babs{f(x)-L} < \epsilon$.
Misalkan $\{ x_n \}_{n=1}^\infty$
suatu barisan dalam $S$ sedemikian sehingga $\lim_{n\to\infty} x_n = \infty$.  Maka terdapat
$N$ sedemikian sehingga
$x_n \geq M$
untuk semua $n \geq N$.
Dengan demikian,
$\babs{f(x_n)-L} < \epsilon$.

Kita membuktikan arah sebaliknya dengan kontraposisi.  Misalkan $f(x)$ tidak
menuju $L$ ketika $x \to \infty$.
Ini berarti bahwa terdapat $\epsilon > 0$
sedemikian sehingga untuk setiap $n \in \N$, terdapat $x \in S$, $x \geq n$, yang kita
namai $x_n$, sedemikian sehingga $\babs{f(x_n)-L} \geq \epsilon$.
Tinjau barisan $\{ x_n \}_{n=1}^\infty$.  Jelas bahwa 
$\bigl\{ f(x_n) \bigr\}_{n=1}^\infty$ tidak konvergen ke $L$.  Tinggal diperhatikan
bahwa $\lim_{n\to\infty} x_n = \infty$ karena $x_n \geq n$ untuk semua $n$.
\end{proof}

Dengan menggunakan lemma tersebut, sekali lagi kita menerjemahkan hasil-hasil mengenai limit
sekuensial menjadi hasil-hasil mengenai limit kontinu ketika $x$ menuju tak hingga.  Artinya,
kita memperoleh analog dari akibat-akibat
dalam \sectionref{subseq:sequentiallimits} secara hampir langsung.  Kita cukup mengizinkan 
titik akumulasi $c$ berupa $\infty$ atau $-\infty$, selain
suatu bilangan real.  Kita menyerahkan
verifikasi pernyataan-pernyataan ini kepada mahasiswa.

\subsection{Limit tak hingga}

Seperti pada barisan, sering kali berguna untuk membedakan beberapa
fungsi divergen dan membicarakan limit yang bernilai tak hingga
seolah-olah limit tersebut ada.

\begin{defn}
\index{limit tak hingga!fungsi}\index{limit!tak hingga}%
Misalkan $f \colon S \to \R$ suatu fungsi dan andaikan
$S$ memiliki $\infty$ sebagai titik akumulasi.
Kita katakan bahwa $f(x)$
\emph{\myindex{divergen ke tak hingga}}
ketika $x$ menuju $\infty$
jika untuk setiap $N \in \R$
terdapat $M \in \R$ sedemikian sehingga
\begin{equation*}
f(x) > N
\end{equation*}
setiap kali $x \in S$ dan $x \geq M$.
Kita tulis
\begin{equation*}
\lim_{x \to \infty} f(x) \coloneqq \infty ,
\end{equation*}
atau kita katakan bahwa $f(x) \to \infty$ ketika $x \to \infty$.
\end{defn}

Definisi serupa dapat dibuat untuk limit ketika $x \to -\infty$
atau ketika $x \to c$ untuk $c$ berhingga.  Definisi serupa juga dapat
dibuat untuk limit yang bernilai $-\infty$.  Perumusan definisi-definisi tersebut
diserahkan sebagai latihan.
Perhatikan bahwa
kadang-kadang digunakan ungkapan yang sedikit tidak tepat (dalam konteks ini),
yaitu
\emph{\myindex{konvergen ke tak hingga}}.
Kita dapat kembali menggunakan limit sekuensial, dan suatu analog dari
\lemmaref{seqflimit:lemma} diserahkan sebagai latihan.

\begin{example}
Mari kita tunjukkan bahwa $\lim\limits_{x \to \infty} \frac{1+x^2}{1+x} = \infty$.

Bukti: Untuk $x \geq 1$, berlaku
\begin{equation*}
\frac{1+x^2}{1+x} \geq 
\frac{x^2}{x+x}  = 
\frac{x}{2} .
\end{equation*}
Untuk sebarang $N \in \R$, pilih $M = \max \{ 2N+1 , 1 \}$.
Jika $x \geq M$, maka $x \geq 1$ dan $\nicefrac{x}{2} > N$.
Jadi,
\begin{equation*}
\frac{1+x^2}{1+x} \geq 
\frac{x}{2} > N .
\end{equation*}
\end{example}

\subsection{Komposisi}

Terakhir, seperti halnya limit pada bilangan berhingga,
limit berperilaku baik terhadap komposisi apabila diberikan
asumsi kekontinuan.

\begin{prop} \label{prop:inflimcompositions}
Andaikan $f \colon A \to B$, $g \colon B \to \R$, $A, B \subset \R$,
$a \in \R \cup \{ -\infty, \infty\}$ merupakan titik akumulasi $A$,
dan $b \in \R \cup \{ -\infty, \infty\}$ merupakan titik akumulasi $B$.
Andaikan
\begin{equation*}
\lim_{x \to a} f(x) = b\qquad \text{dan} \qquad \lim_{y \to b} g(y) = c
\end{equation*}
untuk suatu $c \in \R \cup \{ -\infty, \infty \}$.
Jika $b \in B$, andaikan pula $g(b) = c$.
Maka
\begin{equation*}
\lim_{x \to a} g\bigl(f(x)\bigr) = c .
\end{equation*}
\end{prop}

Pembuktiannya langsung dan diserahkan sebagai latihan.  Kita telah
mengetahui proposisi tersebut ketika $a, b, c \in \R$ (lihat Latihan
\ref{exercise:contlimitcomposition} dan
\ref{exercise:contlimitbadcomposition}).  Sekali lagi, syarat bahwa $g$
kontinu di $b$, jika $b \in B$, memang diperlukan.

\begin{example}
Misalkan $h(x) \coloneqq e^{-x^2+x}$.  Maka
\begin{equation*}
\lim_{x\to \infty} h(x) = 0 .
\end{equation*}

Bukti:
Pernyataan tersebut diperoleh setelah kita mengetahui bahwa
\begin{equation*}
\lim_{x\to \infty} (-x^2+x) = -\infty
\end{equation*}
dan
\begin{equation*}
\lim_{y\to -\infty} e^y = 0 ,
\end{equation*}
yang biasanya dibuktikan ketika fungsi eksponensial didefinisikan.
\end{example}

\subsection{Latihan}

\begin{exercise}
Buktikan \propref{liminfty:unique}.
\end{exercise}

\begin{exercise}
Misalkan $f \colon [1,\infty) \to \R$ suatu fungsi.  Definisikan
$g \colon (0,1] \to \R$ melalui $g(x) \coloneqq f(\nicefrac{1}{x})$.
Dengan menggunakan definisi limit secara langsung,
tunjukkan bahwa $\lim_{x\to 0^+} g(x)$
ada jika dan hanya jika $\lim_{x\to \infty} f(x)$ ada; dalam
hal ini kedua limit tersebut sama.
\end{exercise}

\begin{exercise}
Buktikan \propref{prop:inflimcompositions}.
\end{exercise}

\begin{exercise}
Mari kita jelaskan alasan bagi terminologi tersebut.
Misalkan $f \colon \R \to \R$ suatu fungsi sedemikian sehingga
$\lim_{x \to \infty} f(x) = \infty$ (divergen menuju tak hingga).
Tunjukkan bahwa $f(x)$ divergen (yakni tidak konvergen) ketika $x \to \infty$.
\end{exercise}

\begin{exercise}
Rumuskan definisi limit $f(x)$ menuju $-\infty$ ketika $x \to
\infty$, $x \to -\infty$, dan $x \to c$ untuk $c \in \R$ yang berhingga.
Kemudian nyatakan definisi limit $f(x)$ menuju $\infty$ 
ketika $x \to -\infty$, dan ketika $x \to c$ untuk $c \in \R$ yang berhingga.
\end{exercise}

\begin{exercise}
Andaikan $P(x) \coloneqq x^n + a_{n-1} x^{n-1} + \cdots + a_1 x + a_0$ merupakan \emph{\myindex{polinom monik}}
berderajat $n \geq 1$ (monik berarti koefisien $x^n$ adalah 1).
\begin{enumerate}[a)]
\item
Tunjukkan bahwa jika $n$ genap, maka $\lim\limits_{x\to\infty} P(x) = 
\lim\limits_{x\to-\infty} P(x) = \infty$.
\item
Tunjukkan bahwa jika $n$ ganjil, maka
$\lim\limits_{x\to\infty} P(x) = \infty$ dan
$\lim\limits_{x\to-\infty} P(x) = -\infty$ (lihat latihan sebelumnya).
\end{enumerate}
\end{exercise}

\begin{exercise}
Misalkan $\{ x_n \}_{n=1}^\infty$ suatu barisan.  Tinjau $S \coloneqq \N \subset \R$, dan
$f \colon S \to \R$ yang didefinisikan oleh $f(n) \coloneqq x_n$.  Tunjukkan bahwa
kedua pengertian limit,
\begin{equation*}
\lim_{n\to\infty} x_n \qquad \text{dan} \qquad
\lim_{x\to\infty} f(x) 
\end{equation*}
ekuivalen.  Artinya, tunjukkan bahwa jika salah satunya ada, maka
yang lainnya juga ada, dan dalam hal ini keduanya sama.
\end{exercise}

\begin{exercise}
Perluas \lemmaref{seqflimitinf:lemma} sebagai berikut.
Andaikan $S \subset \R$ memiliki titik akumulasi $c \in \R$, $c = \infty$,
atau $c = -\infty$.  Misalkan $f \colon S \to \R$ suatu fungsi dan andaikan
$L = \infty$ atau $L = -\infty$.  Tunjukkan bahwa
\begin{equation*}
\begin{aligned}
&
\lim_{x\to c} f(x) = L
\\
&
\qquad \text{jika dan hanya jika}
\\
&
\lim_{n\to\infty} f(x_n) = L \enspace \text{untuk setiap barisan }
\{ x_n \}_{n=1}^\infty, \\
& \qquad \text{sedemikian sehingga } \lim_{n\to\infty} x_n = c, \quad
x_n \in S \setminus \{c\} \text{ untuk setiap } n  .
\end{aligned}
\end{equation*}
\end{exercise}

\begin{exercise}
Andaikan $f \colon \R \to \R$ merupakan fungsi periodik dengan periode 2, yakni $f(x + 2) =
f(x)$ untuk setiap $x$.  Definisikan $g \colon \R \to \R$ dengan 
$g(0) \coloneqq f(0)$ dan, untuk $x \neq 0$, dengan
\begin{equation*}
g(x) \coloneqq f\left(\frac{\sqrt{x^2+1}-1}{x}\right)
\end{equation*}
\begin{enumerate}[a)]
\item
Tentukan fungsi $\varphi \colon (-1,1) \to \R$ sedemikian sehingga
$g\bigl(\varphi(t)\bigr) = f(t)$, yakni $\varphi^{-1}(x) = 
\frac{\sqrt{x^2+1}-1}{x}$ untuk $x \neq 0$.
\item
Tunjukkan bahwa $f$ kontinu jika dan hanya jika $g$ kontinu dan
\begin{equation*}
\lim_{x \to \infty} g(x) = 
\lim_{x \to -\infty} g(x) = 
f(1) = f(-1) .
\end{equation*}
\end{enumerate}
\end{exercise}

%%%%%%%%%%%%%%%%%%%%%%%%%%%%%%%%%%%%%%%%%%%%%%%%%%%%%%%%%%%%%%%%%%%%%%%%%%%%%%

\sectionnewpage
\section{Fungsi monoton dan kekontinuan}
\label{sec:monotonefunc}

%mbxINTROSUBSECTION

\sectionnotes{1 pertemuan (opsional, dapat dihilangkan dengan aman kecuali
\sectionref{sec:ift} juga dibahas, memerlukan \sectionref{sec:limitatinf})}

\begin{defn}
Misalkan $S \subset \R$.
Kita katakan bahwa $f \colon S \to \R$ \emph{\myindex{monoton naik}}
(berturut-turut, \emph{\myindex{naik tegas}}) jika untuk $x,y \in S$ dengan
$x < y$ berlaku $f(x) \leq f(y)$ (berturut-turut, $f(x) < f(y)$).
Kita definisikan
\emph{\myindex{monoton turun}} dan
\emph{\myindex{turun tegas}} dengan cara yang sama, yaitu dengan membalik
ketaksamaan untuk $f$.

Jika suatu fungsi monoton naik atau monoton turun, kita katakan fungsi tersebut
\emph{monoton}\index{fungsi monoton}.  Jika fungsi tersebut
naik tegas atau turun tegas, kita katakan fungsi tersebut
\emph{monoton tegas}\index{fungsi monoton tegas}.
\end{defn}

Kadang-kadang istilah \emph{\myindex{tak menurun}}
(berturut-turut, \emph{\myindex{tak menaik}}) digunakan
untuk fungsi monoton naik (berturut-turut, monoton turun) guna menekankan bahwa fungsi itu tidak
harus naik tegas (berturut-turut, turun tegas).

Jika $f$ monoton naik, maka $-f$ monoton turun, dan sebaliknya.  Oleh karena itu,
banyak hasil mengenai fungsi monoton cukup dibuktikan, misalnya, untuk fungsi
monoton naik, lalu hasil untuk fungsi monoton turun diperoleh dengan mudah.

\subsection{Kekontinuan fungsi monoton}

Limit satu sisi untuk fungsi monoton dihitung
dengan menghitung infimum dan supremum.

\begin{prop} \label{prop:monotlimits}
Misalkan $S \subset \R$, $c \in \R$,
$f \colon S \to \R$ monoton naik,
dan
$g \colon S \to \R$ monoton turun.
Jika $c$ merupakan titik akumulasi dari $S \cap (-\infty,c)$, maka
\begin{equation*}
\lim_{x \to c^-} f(x) = \sup \{ f(x) : x < c, x \in S \}
%mbxSTARTIGNORE
\qquad \text{dan} \qquad
%mbxENDIGNORE
%mbxlatex \quad \text{and} \quad
\lim_{x \to c^-} g(x) = \inf \{ g(x) : x < c, x \in S \} .
\end{equation*}
Jika $c$ merupakan titik akumulasi dari $S \cap (c,\infty)$, maka
\begin{equation*}
\lim_{x \to c^+} f(x) = \inf \{ f(x) : x > c, x \in S \}
%mbxSTARTIGNORE
\qquad \text{dan} \qquad
%mbxENDIGNORE
%mbxlatex \quad \text{and} \quad
\lim_{x \to c^+} g(x) = \sup \{ g(x) : x > c, x \in S \} .
\end{equation*}
Jika $\infty$ merupakan titik akumulasi dari $S$, maka
\begin{equation*}
\lim_{x \to \infty} f(x) = \sup \{ f(x) : x \in S \}
%mbxSTARTIGNORE
\qquad \text{dan} \qquad
%mbxENDIGNORE
%mbxlatex \quad \text{and} \quad
\lim_{x \to \infty} g(x) = \inf \{ g(x) : x \in S \} .
\end{equation*}
Jika $-\infty$ merupakan titik akumulasi dari $S$, maka
\begin{equation*}
\lim_{x \to -\infty} f(x) = \inf \{ f(x) : x \in S \}
%mbxSTARTIGNORE
\qquad \text{dan} \qquad
%mbxENDIGNORE
%mbxlatex \quad \text{and} \quad
\lim_{x \to -\infty} g(x) = \sup \{ g(x) : x \in S \} .
\end{equation*}
\end{prop}

Dengan kata lain, semua limit satu sisi ada setiap kali limit tersebut
terdefinisi.  Oleh karena itu, untuk fungsi monoton, ketika kita mengatakan
bahwa limit kiri $x \to c^-$
ada, yang kita maksud ialah bahwa $c$ merupakan titik akumulasi dari $S \cap (-\infty,c)$,
dan demikian pula untuk limit kanan.

\begin{proof}
Misalkan $f$ monoton naik, dan kita akan menunjukkan kesamaan
pertama.  Bagian lain dari bukti ini sangat serupa dan diserahkan sebagai
latihan.

Misalkan $a \coloneqq \sup \{ f(x) : x < c, x \in S \}$.  Jika $a = \infty$,
maka untuk sebarang $M \in \R$, terdapat $x_M \in S$, $x_M < c$, sedemikian sehingga $f(x_M) > M$. 
Karena $f$ monoton naik, $f(x) \geq f(x_M) >  M$ untuk setiap $x \in S$ dengan $x > x_M$.
Ambil $\delta \coloneqq c-x_M > 0$ agar definisi limit yang
bernilai tak hingga terpenuhi.

Selanjutnya, misalkan $a < \infty$.
Ambil sebarang $\epsilon > 0$.  Karena $a$ merupakan supremum dan
$S \cap (-\infty,c)$ tidak kosong, $a \in \R$ dan
terdapat
$x_\epsilon \in S$,
$x_\epsilon < c$,
sedemikian sehingga $f(x_\epsilon) > a-\epsilon$.  Karena $f$ monoton naik,
jika $x \in S$ dan $x_\epsilon < x < c$, maka
$a-\epsilon < f(x_\epsilon) \leq f(x) \leq a$.  Misalkan
$\delta \coloneqq c-x_\epsilon$.  Maka untuk $x \in S \cap (-\infty,c)$
dengan $\sabs{x-c} < \delta$,
berlaku $\babs{f(x)-a} < \epsilon$.
\end{proof}

Misalkan $f \colon S \to \R$ monoton naik, $c \in S$, dan
kedua limit satu sisinya ada.
Karena $f(x) \leq f(c) \leq f(y)$
setiap kali $x < c < y$, dengan mengambil limit kita memperoleh
\begin{equation*}
\lim_{x \to c^-} f(x) \leq f(c) \leq \lim_{x \to c^+} f(x) .
\end{equation*}
Maka $f$ kontinu di $c$ jika dan hanya jika kedua limit tersebut sama
satu sama lain (dan karenanya sama dengan $f(c)$).  Lihat juga
\propref{prop:onesidedlimits}.
Lihat \figureref{fig:figinccont} untuk memperoleh gambaran mengenai seperti apa
suatu diskontinuitas.


\begin{cor} \label{cor:continterval}
Jika $I \subset \R$ merupakan interval dan $f \colon I \to \R$ 
monoton dan tidak konstan, maka $f(I)$ merupakan interval jika dan hanya jika $f$
kontinu.
\end{cor}

Asumsi bahwa $f$ tidak konstan digunakan untuk menghindari persoalan teknis
bahwa $f(I)$ merupakan satu titik: $f(I)$ merupakan satu
titik jika dan hanya jika $f$ konstan.  Fungsi konstan bersifat 
kontinu.

\begin{proof}
Tanpa mengurangi keumuman, misalkan $f$ monoton naik.

Pertama, misalkan $f$ kontinu.  Ambil dua titik
$f(x_1), f(x_2)$ di $f(I)$ dan misalkan
$f(x_1) < f(x_2)$.
Karena $f$ monoton naik, $x_1 < x_2$.  Berdasarkan
\hyperref[IVT:thm]{teorema nilai antara},
untuk sebarang $y$ dengan $f(x_1) < y < f(x_2)$, kita menemukan
$c \in (x_1,x_2) \subset I$ sedemikian sehingga $f(c) = y$, jadi $y \in f(I)$. 
Dengan demikian, $f(I)$ merupakan interval.

Mari kita buktikan arah sebaliknya dengan kontraposisi.
Misalkan $f$ tidak kontinu di $c \in I$,
dan $c$ bukan titik ujung dari $I$.
Misalkan
\begin{equation*}
%mbxlatex \begin{aligned}
a
%mbxlatex &
\coloneqq \lim_{x \to c^-} f(x) = \sup \bigl\{ f(x) : x \in I, x < c \bigr\} ,
%mbxSTARTIGNORE
\qquad
%mbxENDIGNORE
%mbxlatex \\
b
%mbxlatex &
\coloneqq \lim_{x \to c^+} f(x) = \inf \bigl\{ f(x) : x \in I, x > c \bigr\} .
%mbxlatex \end{aligned}
\end{equation*}
Karena $c$ merupakan titik diskontinuitas, $a < b$.
Jika $x < c$, maka $f(x) \leq a$, dan
jika $x > c$, maka $f(x) \geq b$.  Oleh karena itu,
tidak ada titik
di $(a,b) \setminus \bigl\{ f(c) \bigr\}$ yang termasuk dalam $f(I)$.
Terdapat $x_1 \in I$ dengan $x_1 < c$, sehingga
$f(x_1) \leq a$, dan terdapat $x_2 \in I$ dengan $x_2 > c$,
sehingga $f(x_2) \geq b$.  Baik $f(x_1)$ maupun $f(x_2)$ termasuk dalam $f(I)$,
tetapi terdapat titik-titik di antaranya yang tidak termasuk dalam $f(I)$.
Jadi, $f(I)$ bukan merupakan interval.  Lihat \figureref{fig:figinccont}.

Jika $c \in I$ merupakan titik ujung, buktinya serupa dan diserahkan sebagai latihan.
\end{proof}

\begin{myfigureht}
\myincludepdft{figinccont}{%
Diagram grafik fungsi dengan diskontinuitas loncatan beserta pengaruhnya
terhadap citra fungsi tersebut.
Fungsi ini didefinisikan pada interval bertanda I.  Titik x subskrip 1, c, dan
x subskrip 2 ditandai di dalam I dan tersusun dalam urutan tersebut dari kiri ke kanan.
Fungsi naik secara kontinu sampai x sama dengan c,
tempat fungsi mendekati a; yaitu, limit f dari x ketika x mendekati
c dari kiri adalah a.  Nilai f di c adalah f dari c, yang lebih besar daripada a
dan lebih kecil daripada b, dengan b merupakan limit f dari x ketika x mendekati c dari
kanan.  Kemudian fungsi terus naik secara kontinu hingga ujung
interval I.  Citra f terbagi menjadi tiga bagian.
Bagian terendah merupakan
interval yang batas atasnya a, dan interval ini memuat
f dari x subskrip 1.
Bagian berikutnya hanyalah titik f dari c.
Bagian ketiga, yang tertinggi, merupakan interval yang berawal di b,
dan interval ini memuat f dari x subskrip 2.}
\caption{Diskontinuitas fungsi monoton naik $f \colon I \to \R$ di
$c$.\label{fig:figinccont}}
\end{myfigureht}

Salah satu sifat mencolok dari fungsi monoton adalah bahwa fungsi tersebut tidak mungkin memiliki
terlalu banyak diskontinuitas.

\begin{cor} \label{cor:monotcountcont}
Misalkan $I \subset \R$ merupakan interval dan
$f \colon I \to \R$ monoton.  Maka $f$ memiliki himpunan
titik diskontinuitas yang terhitung.
\end{cor}

\begin{proof}
Misalkan $E \subset I$ merupakan himpunan semua titik diskontinuitas
yang bukan titik ujung dari $I$.  Karena hanya ada
dua titik ujung, cukup ditunjukkan bahwa $E$ terhitung.
Tanpa mengurangi keumuman, misalkan $f$ monoton naik.
Kita akan mendefinisikan suatu injeksi $h \colon E \to \Q$.
Untuk setiap $c \in E$,
kedua limit satu sisi dari $f$ ada karena $c$ bukan titik ujung.
Misalkan
\begin{equation*}
%mbxlatex \begin{aligned}
a
%mbxlatex &
\coloneqq \lim_{x \to c^-} f(x) = \sup \bigl\{ f(x) : x \in I, x < c \bigr\} ,
%mbxSTARTIGNORE
\qquad
%mbxENDIGNORE
%mbxlatex \\
b
%mbxlatex &
\coloneqq \lim_{x \to c^+} f(x) = \inf \bigl\{ f(x) : x \in I, x > c \bigr\} .
%mbxlatex \end{aligned}
\end{equation*}
Karena $c$ merupakan titik diskontinuitas, $a < b$.  
Terdapat bilangan rasional $q \in (a,b)$, jadi tetapkan $h(c) \coloneqq q$.
Misalkan $d \in E$ merupakan titik diskontinuitas lain.
Jika $d > c$, terdapat
$x \in I$ dengan $c < x < d$, sehingga $\lim_{x \to d^-} f(x) \geq b$.
Oleh karena itu, bilangan rasional yang kita pilih untuk $h(d)$ berbeda dari $q$,
sebab $q=h(c) < b$ dan $h(d) > b$.
Demikian pula jika $d < c$.  Setelah membuat pilihan seperti itu untuk
setiap unsur $E$, kita memperoleh fungsi 
satu-ke-satu (injektif) ke dalam $\Q$.  Oleh karena itu, $E$ terhitung.
\end{proof}

\begin{example} \label{example:countdiscont}
Nyatakan bilangan bulat terbesar yang kurang dari atau sama dengan $x$
dengan $\lfloor x \rfloor$.
Definisikan $f \colon [0,1] \to \R$ dengan
\begin{equation*}
f(x) \coloneqq
x +
\sum_{n=0}^{\lfloor 1/(1-x) \rfloor}
2^{-n} ,
\end{equation*}
untuk $x < 1$ dan $f(1) \coloneqq 3$.
Sebagai latihan, tunjukkan bahwa $f$ naik tegas, terbatas, dan
memiliki diskontinuitas di setiap titik $1-\nicefrac{1}{k}$ untuk $k \in \N$ dengan $k \geq 2$.  Secara khusus,
himpunan titik diskontinuitas fungsi tersebut terhitung, tetapi fungsi itu terbatas dan
didefinisikan pada interval tertutup terbatas.  Lihat \figureref{fig:countdiscont}.
\begin{myfigureht}
\myincludegraphics{increasing-discont-fig}{%
Grafik fungsi pada interval dari 0 hingga 1
yang tersusun atas ruas-ruas garis lurus yang miring ke atas.  Terdapat
ruas pertama yang melintasi separuh grafik, berawal pada 
y sama dengan 1,5 dan berakhir pada y sama dengan 2.  Kemudian terdapat
diskontinuitas loncatan ke atas dan ruas garis lain yang jauh lebih pendek,
yang juga miring ke atas.  Rangkaian loncatan ke atas yang makin kecil dan
ruas garis miring ke atas yang makin pendek ini terus berlanjut hingga
pada x sama dengan 1 kita mencapai y sama dengan 3.}
\caption{Fungsi naik tegas pada $[0,1]$ dengan himpunan titik diskontinuitas
yang terhitung.\label{fig:countdiscont}}
\end{myfigureht}

Dengan cara serupa, kita dapat menemukan contoh fungsi monoton yang diskontinu pada himpunan rapat
seperti bilangan rasional.  Lihat latihan.
\end{example}


\subsection{Kekontinuan fungsi invers}


Fungsi monoton tegas $f$ bersifat satu-satu (injektif).
Untuk melihat fakta ini,
perhatikan bahwa jika $x \neq y$, kita dapat mengasumsikan $x < y$.  Kita memperoleh $f(x) <
f(y)$ jika $f$ naik tegas, atau $f(x) > f(y)$ jika $f$ turun
tegas, sehingga $f(x) \neq f(y)$.
Oleh karena itu, $f$ harus memiliki invers $f^{-1}$ yang terdefinisi pada daerah hasilnya.

\begin{prop} \label{prop:invcont}
Jika $I \subset \R$ adalah suatu interval dan $f \colon I \to \R$ monoton
tegas, maka invers $f^{-1} \colon f(I) \to I$ kontinu.
\end{prop}

\begin{proof}
Misalkan $f$ naik tegas.  Buktinya hampir
sama untuk fungsi yang turun tegas.
Karena $f$ naik tegas, demikian pula $f^{-1}$.  Artinya, jika $f(x) <
f(y)$, maka harus berlaku $x < y$ dan oleh karena itu
$f^{-1}\bigl(f(x)\bigr) < f^{-1}\bigl(f(y)\bigr)$.

Ambil $c \in f(I)$.
Jika $c$ bukan titik akumulasi dari $f(I)$, maka $f^{-1}$ kontinu di $c$
secara otomatis.  Jadi, misalkan $c$ adalah titik akumulasi dari $f(I)$.
Misalkan kedua limit satu sisi berikut ada, yakni keduanya
terdefinisi:
\begin{align*}
x_0 & \coloneqq \lim_{y \to c^-} f^{-1}(y) =
\sup \bigl\{ f^{-1}(y) : y < c, y \in f(I) \bigr\}
=
\sup \bigl\{ x \in I : f(x) < c \bigr\} , \\
x_1 & \coloneqq \lim_{y \to c^+} f^{-1}(y) =
\inf \bigl\{ f^{-1}(y) : y > c, y \in f(I) \bigr\}
=
\inf \bigl\{ x \in I : f(x) > c \bigr\} .
\end{align*}
Kita memiliki $x_0 \leq x_1$ karena $f^{-1}$ monoton naik.
Untuk setiap $x \in I$ dengan $x > x_0$, berlaku $f(x) \geq c$.  Karena $f$ naik tegas,
harus berlaku $f(x) > c$ untuk setiap $x \in I$ dengan $x > x_0$.  Oleh karena itu,
\begin{equation*}
\{ x \in I : x > x_0 \} \subset \bigl\{ x \in I : f(x) > c \bigr\}.
\end{equation*}
Infimum himpunan di sebelah kiri adalah $x_0$, dan infimum himpunan di sebelah kanan
adalah $x_1$, sehingga kita memperoleh $x_0 \geq x_1$.
Jadi, $x_1 = x_0$, dan $f^{-1}$ kontinu di $c$.

Jika salah satu limit satu sisi tidak ada (yakni, ketika pengambilan limit dari
sisi tersebut tidak terdefinisi), argumennya serupa
dan diserahkan sebagai latihan.
\end{proof}

\begin{example}
Proposisi tersebut tidak mensyaratkan bahwa $f$ sendiri kontinu.  Misalkan
$f \colon \R \to \R$ didefinisikan oleh
\begin{equation*}
f(x) \coloneqq
\begin{cases}
x & \text{jika } x < 0, \\
x+1 & \text{jika } x \geq 0. \\
\end{cases}
\end{equation*}
Fungsi $f$ tidak kontinu di $0$.
Citra dari $I = \R$ adalah himpunan 
$(-\infty,0)\cup [1,\infty)$, yang bukan suatu interval.
Dengan demikian, $f^{-1} \colon (-\infty,0)\cup [1,\infty)
\to \R$ dapat ditulis sebagai
\begin{equation*}
f^{-1}(y) =
\begin{cases}
y & \text{jika } y < 0, \\
y-1 & \text{jika } y \geq 1. 
\end{cases}
\end{equation*}
Tidak sulit untuk melihat bahwa $f^{-1}$ adalah fungsi kontinu.  Lihat
\figureref{invcontfig} untuk grafik-grafiknya.
\begin{myfigureht}
\myincludepdft{invcontfigAB}{%
Dua grafik.  Di sebelah kiri, sebuah garis diagonal menanjak yang membentang dari
kiri bawah ke titik asal (tidak termasuk titik asal),
lalu meloncat ke atas dan kemudian terus bergerak diagonal ke atas.
Pada grafik di sebelah kanan, garis pertama yang membentuk
grafik itu sama, tetapi kemudian terdapat sebuah interval pendek yang tidak dilalui
grafik dan garisnya dimulai kembali pada sumbu horizontal,
lalu bergerak diagonal ke atas lagi.}
\caption{Grafik $f$ di sebelah kiri dan $f^{-1}$ di sebelah kanan.\label{invcontfig}}
\end{myfigureht}
\end{example}

Perhatikan apa yang terjadi pada proposisi tersebut jika $f(I)$ merupakan suatu interval.
Dalam hal itu, kita cukup
menerapkan \corref{cor:continterval} pada $f$ dan $f^{-1}$.  Artinya, jika
$f \colon I \to J$ adalah fungsi monoton tegas yang surjektif dan $I$ serta $J$ merupakan interval,
maka $f$ dan $f^{-1}$ keduanya kontinu.  Lebih lanjut, $f(I)$ merupakan suatu
interval tepat ketika $f$ kontinu.

\subsection{Latihan}

\begin{exercise}
Misalkan $f \colon [0,1] \to \R$ monoton.  Buktikan bahwa $f$ terbatas.
\end{exercise}

\begin{exercise}
Lengkapi bukti \propref{prop:monotlimits}.
Petunjuk: Anda dapat mengurangi pekerjaan menjadi separuhnya dengan memperhatikan bahwa jika $g$ monoton turun,
maka $-g$ monoton naik.
\end{exercise}

\begin{exercise}
Lengkapi bukti \corref{cor:continterval}.
\end{exercise}

\begin{exercise}
Buktikan pernyataan-pernyataan dalam \exampleref{example:countdiscont}.
\end{exercise}

\begin{exercise}
Lengkapi bukti \propref{prop:invcont}.
\end{exercise}

\begin{samepage}
\begin{exercise}
Misalkan $S \subset \R$, dan $f \colon S \to \R$ merupakan fungsi
monoton naik.  Buktikan:
\begin{enumerate}[a)]
\item
Jika $c$ merupakan titik akumulasi
dari $S \cap (c,\infty)$, maka
$\lim\limits_{x\to c^+} f(x) < \infty$.
\item
Jika $c$ merupakan titik akumulasi dari $S \cap (-\infty,c)$
dan $\lim\limits_{x\to c^-} f(x) = \infty$, maka
$S \subset (-\infty,c)$.
\end{enumerate}
\end{exercise}
\end{samepage}

\begin{exercise}
Misalkan $I \subset \R$ merupakan suatu interval dan $f \colon I \to \R$ suatu fungsi.
Andaikan untuk setiap $c \in I$, terdapat $a, b \in \R$ dengan
$a > 0$ sedemikian sehingga $f(x) \geq a x + b$ untuk setiap $x \in I$
dan $f(c) = a c + b$.  Tunjukkan bahwa $f$ naik tegas.
\end{exercise}

\begin{exercise}
Misalkan $I$ dan $J$ merupakan interval dan
$f \colon I \to J$ merupakan fungsi kontinu dan bijektif (satu-satu dan surjektif).
Tunjukkan bahwa $f$ monoton tegas.
\end{exercise}

\begin{exercise}
Tinjau suatu fungsi monoton $f \colon I \to \R$ pada interval $I$.  Buktikan bahwa terdapat
fungsi $g \colon I \to \R$ sedemikian sehingga
$\lim\limits_{x \to c^-} g(x) = g(c)$ untuk setiap $c$ dalam $I$ kecuali
titik ujung yang lebih kecil (kiri) dari $I$, dan sedemikian sehingga
$g(x) = f(x)$ untuk setiap $x \in I$ di luar suatu himpunan terhitung.
\end{exercise}

\begin{exercise}
\leavevmode
\begin{enumerate}[a)]
\item
Misalkan $S \subset \R$ suatu himpunan bagian.  Jika $f \colon S \to \R$ monoton naik dan
terbatas,
maka tunjukkan bahwa terdapat fungsi monoton naik $F \colon \R \to \R$
sedemikian sehingga $f(x) = F(x)$ untuk setiap $x \in S$.
\item
Temukan contoh fungsi $f \colon S \to \R$ yang terbatas dan naik tegas sedemikian sehingga
fungsi monoton naik $F$ seperti di atas tidak dapat naik tegas, bagaimanapun kita
memilihnya.
\end{enumerate}
\end{exercise}

\begin{exercise}[Menantang] \label{exercise:increasingfuncdiscatQ}
Temukan contoh fungsi monoton naik $f \colon [0,1] \to \R$
yang memiliki diskontinuitas pada setiap bilangan rasional.  Kemudian tunjukkan bahwa citra
$f\bigl([0,1]\bigr)$ tidak memuat interval apa pun.  Petunjuk: Enumerasikan
bilangan-bilangan rasional dan definisikan
fungsi tersebut dengan suatu deret.
\end{exercise}

\begin{exercise}
Misalkan $I$ merupakan suatu interval dan $f \colon I \to \R$ monoton.
Tunjukkan bahwa $\R \setminus f(I)$ merupakan gabungan terhitung dari interval-interval yang saling lepas.
\end{exercise}

\begin{exercise}
Misalkan $f \colon [0,1] \to (0,1)$ monoton naik.  Tunjukkan bahwa untuk setiap
$\epsilon > 0$, terdapat
fungsi naik tegas $g \colon [0,1] \to (0,1)$ sedemikian sehingga
$g(0) = f(0)$, $f(x) \leq g(x)$ untuk setiap $x$, dan $g(1)-f(1) < \epsilon$.
\end{exercise}

\begin{exercise}
Buktikan bahwa \myindex{fungsi Dirichlet}
$f \colon [0,1] \to\R$, yang didefinisikan oleh $f(x) \coloneqq 1$
jika $x$ rasional dan $f(x) \coloneqq 0$ selainnya, tidak dapat ditulis sebagai
selisih dua fungsi monoton naik.  Artinya, tidak terdapat fungsi monoton naik
$g$ dan $h$ sedemikian sehingga $f(x) = g(x) - h(x)$.
\end{exercise}

\begin{exercise}
Misalkan $f \colon (a,b) \to (c,d)$ merupakan fungsi naik tegas
dan surjektif.  Buktikan bahwa terdapat $g \colon (a,b) \to (c,d)$
yang juga naik tegas dan surjektif, serta $g(x) < f(x)$ untuk setiap $x \in
(a,b)$.
\end{exercise}
